\documentclass[11pt,reqno]{amsart}
\usepackage[T1]{fontenc}
\usepackage{lmodern,amsmath,amssymb,mathtools,microtype,booktabs,array,longtable}
\usepackage[margin=1in]{geometry}
\usepackage{graphicx}
\usepackage{xr-hyper}
\usepackage[hidelinks]{hyperref}
\ifdefined\pdfinfoomitdate\pdfinfoomitdate=1\fi
\ifdefined\pdftrailerid\pdftrailerid{}\fi
\ifdefined\pdfsuppressptexinfo\pdfsuppressptexinfo=15\fi
\numberwithin{equation}{section}
\newtheorem{theorem}{Theorem}[section]
\newtheorem{proposition}[theorem]{Proposition}
\newtheorem{lemma}[theorem]{Lemma}
\newtheorem{corollary}[theorem]{Corollary}
\theoremstyle{definition}\newtheorem{definition}[theorem]{Definition}
\newtheorem{example}[theorem]{Example}
\theoremstyle{remark}\newtheorem{remark}[theorem]{Remark}
\newcommand{\E}{\mathbb E}\newcommand{\Prb}{\mathbb P}
\newcommand{\R}{\mathbb R}\newcommand{\one}{\mathbf1}
\newcommand{\Q}{\mathcal Q}\newcommand{\TV}{\mathrm{TV}}
\newcommand{\cp}{\mathrm{cp}}\newcommand{\ev}{\mathrm{ev}}\newcommand{\st}{\mathrm{st}}
\DeclareMathOperator{\tr}{tr}\DeclareMathOperator{\dist}{dist}
\DeclareMathOperator{\law}{Law}\DeclareMathOperator{\rank}{rank}
\DeclareMathOperator{\supp}{supp}
\emergencystretch=2em

\externaldocument[][nocite]{supplement-refs}[General_Theta_Foundations_I_R73_Supplement.pdf]
\hypersetup{pdftitle={General Theta Foundations I: Acquired Geometry and Causal Resource Transfer},pdfauthor={Qian Qi}}
\title[Acquired geometry and causal resource transfer]{General Theta Foundations I:\\Acquired Geometry and Causal Resource Transfer}
\author{Qian Qi}\date{October 2026}
\subjclass[2020]{Primary 60J10, 68Q45; Secondary 20C15, 60G40, 62C20, 81P15, 90C22, 94A29}
\keywords{causal experiment, observable response, stochastic realization, compact group, minimal memory, algebraic synthesis, quantum instrument}
\begin{document}
\begin{abstract}
We study the finite classical memory required to reproduce all commanded
terminal laws of a prepared causal experiment. For a finite-dimensional
compactly reversible response, every finite stochastic realization admits,
without increasing its number of live labels, a finite permutation core
with at most one transient initializer. Thus continuing transition
randomness can be replaced by one counted initial randomization. This
gives an exact variational formula for the least memory over finite actions
of the observable group, sharp factorial group bounds, and rigidity at
linear-rank equality. Two commands give algebraic families with exactly
$k+1$ stochastic labels and $k!$ deterministic-transition labels in every
dimension $k\ge3$. A qubit example shows that bounded decoder constraints
can strictly strengthen both rank and permutation-module obstructions.
For algebraic observable data, finite realizability is decidable and a
smallest controller can be synthesized exactly. Finite word tests give
positive, algebraically certifiable memory gaps and a computable interval
of unchanged exact memory under damping toward the group-averaged response.
The proof separates classical stochastic matrix-group structure from the
compression needed to preserve the observable law and its prescribed initial
state. Complete periodic, spectral, leakage, stopping, minimax, reset, rank
and geometric results remain available with their individual hypotheses
and resource accounts.
\end{abstract}
\maketitle
\section{Introduction}\label{sec:current-overview}\label{sec:introduction}
A prepared experiment assigns a terminal probability law to each finite
word of commands. Even when its protected state has finite dimension,
an exact classical realization of all these laws can require infinitely
many retained labels. When a finite realization does exist, three sizes
must still be distinguished: the linear response dimension, the number
of distinct deterministic residuals, and the smallest stochastic memory.
The last can be strictly between the first two. We determine the form of
that stochastic advantage for reversible observable dynamics.

Let $b:A^*\to\Delta(C)$ be a probability response on finite command and
terminal alphabets. Its minimal linear space is spanned by the residual
functions $b^w(z)=b(wz)$. We call $b$ compactly reversible when its command
shifts are invertible on this space and generate a group with compact
closure. Prepared finite-dimensional unitary commands always have this
property after directions invisible to every future commanded terminal
test have been removed. The condition concerns observable responses;
the full physical unitary group may be larger.

The finite realization criterion in Theorem~\ref{thm:r72-reversible}
states that this compact observable group must be finite. The main new
result is more precise. Theorem~\ref{thm:r73-normal-form} replaces any
$M$-label stochastic realization by a permutation core and at most one
transient initializer, using at most $M$ labels in total. Its proof does
not merely apply a matrix-group theorem to the supplied transition
matrices: those matrices need not form a group, and their inaccessible
hidden responses need not belong to the minimal observable space. We
average words acting identically on the observable response, pass to the
closed classes of that finite chain, and prove that the compressed
command matrices factor through the observable group. The prescribed
initial state determines whether a further initializer is necessary.

The resulting variational formula, Theorem~\ref{thm:r73-memory-formula},
minimizes over finite permutation actions and bounded probability
decoders. It is an exact optimization over all real stochastic
competitors. In particular $|G_b|\le M!$, and a strict stochastic saving
over the deterministic residual count requires
$|G_b|\le(M-1)!$. Both bounds are sharp. A two-command family on $k$
coordinates has exactly $k+1$ stochastic live labels and $k!$
deterministic-transition labels for every $k\ge3$. The family has
rational preparations and readouts and admits a diagonal physical
realization. It exhibits a stochastic implementation advantage; it is
not presented as a quantum advantage over a classical experiment.

Probability constraints in the factorization carry information beyond the
group and its linear representation. Proposition~\ref{prop:r73-decoder-obstruction}
gives two responses with the same observable group $S_3$, linear dimension
$3$ and six deterministic residuals, but with stochastic minima $4$ and $6$.
The latter has an explicit protected-qubit realization. Its endpoint laws
force incompatible decoder supports in every action with at most five
points. Thus neither residual counting nor permutation-module dimension
alone determines the stochastic resource.

Theorem~\ref{thm:r73-algebraic-synthesis} makes the classification
effective for supplied algebraic observable data. A finite-group order
bound gives a terminating test of existence of any finite realization.
The normal form then reduces exact minimal synthesis to finitely many
finite-group actions and bounded bilinear probability equations.
The returned controller has algebraic entries and a counted initial
label. Theorem~\ref{thm:r73-finite-tests} and the ensuing gap theorem
give finite word certificates against arbitrary bounded-size stochastic
competitors. Applied to damping toward the group-averaged response, they
upgrade qualitative compactness to a computable positive interval on
which the exact minimal memory is unchanged. This interval is not
claimed to have the sharp first transition point.

\subsection*{Classical structure and the present assertion}
Finite residual equivalence for deterministic output machines is
classical \cite{RabinScott59}. Positive realization, observable operator
models and finite linear representations also have an extensive history
\cite{Heller65,CarlylePaz71,Jaeger00,ThonJaeger15,MonrasWinter14}.
Schwarz classified stochastic matrix groups, and Flor gave an elementary
description of nonnegative matrix groups, including the bounded-group
permutation structure and Brown's finiteness theorem
\cite{Schwarz64,Flor69,Brown64}. These results are prior mathematics.
Group images of completely simple semigroups also have a classical
structure theory \cite{Schwarz62}.
The issue addressed here is how an arbitrary stochastic realization of a
reversible \emph{observable} response can be compressed to such a group
while preserving a pure initial state, its terminal law, and the exact
label budget. The initializer-sensitive formula and its constructive
consequences are the statements for which a contribution is asserted.

Finite matrix-group recognition, the congruence bounds of Minkowski, and
real-algebraic automaton minimization remain credited to their existing
theories \cite{Serre07,BabaiBealsRockmore93,KieferWachter14}. The
finite-word equivalence argument is included as a classical lemma
\cite{Tzeng92}. Likewise, finite-dimensional quantum processes without
finite classical realizations and noise-dependent memory lower bounds
predate this paper \cite{Fanizza24}. Neither a general quantum/classical
separation nor those algorithmic ingredients are claimed as discoveries
of the present work. The statements and proofs below permit a direct
comparison with that literature; an exhaustive priority or editorial
verdict is not a mathematical consequence of them.

\subsection*{Relation to the causal resource theory}
The exact periodic partition law, its one-label leakage increment, and
the single-frequency and noncommuting-pair arithmetic envelopes remain
complete current results. Their proofs are in the paired supplement
except for the all-word reversible criterion and common-command
construction retained in the main article. The new normal form places
finite cyclic cases within a general finite-group realization problem,
while strictly damped irrational modes remain outside the reversible
class and retain their distinct spectral arguments.

The main article also retains the full actual-history causal comparison
and common paid construction. Its lower witness is formed on the
competing experiment's actual pre-call histories; it does not assume a
terminating reference run on another experiment. The supplement retains
performed-law compactness, boundary primal and dual attainment, robust
tied optimality, normalized stopping rigidity, instrument realization,
finite-width transfer, reset control, rank-constrained discrimination,
spectral geometry and composition. The independently compiled companion
retains absolute calibration and nonintervening directional sensing.
These results keep their individual hypotheses. Exact response
simulation is a component of this resource theory and is not a substitute
for a general classification of optimal feedback experiments.

\section{Prepared kernels, acquired task laws and observers}\label{sec:data}
All hidden-state, action and report spaces are standard Borel. A fixed unknown index $\theta\in\Theta$ specifies an actual preparation and positive normalized causal kernels. Positivity means nonnegative measures, not an everywhere strictly positive density. Dependence on the entire physical past may be included in the hidden state. A legal policy sees past visible reports and its admitted randomization, but not $\theta$ or a future audit. Iteration of these kernels defines its finite-prefix laws; the consistent countable path law is their kernel extension.

Each call has a cost $c_i\ge c_*>0$. The decision to make the next call is predictable. In the geometric route, the declared productive-call flag is also predictable; in the primitive-instrument route, productive reports are specified outcomes and can be random. Commitment is permitted only after a productive report, and occurs exactly once almost surely in each model. At that time a specified future audit $Z\in H$ is scored, where $H\subset\R^k$ is nonempty, compact and convex. The prediction is committed before that audit is observed. A prospective joint report--audit kernel is defined at every productive call, including rejected calls. The stopping coin is independent of the audit conditional on the entire current visible report and the past. Write $L=\operatorname{diam}(H)^2$; the case $L=0$ is trivial whenever a legal rule exists.

The core cost $C_\tau$ includes initial preparation charged to the core, failed calls and all productive calls through commitment. A separately executed final audit must be added to the raw budget. This convention is fixed before taking a risk infimum. Physical time, stored program constants and temporary computation are not identified with call count.

\subsection*{Three resource interfaces}
A checkpoint procedure may use its whole precommit history, but must output one element of a fixed dictionary of at most $M$ values. This counts its output width only; it is not a claim of $\log_2 M$ total working bits.

An autonomous observer has an alphabet $S$, $|S|\le M$, and one fixed parameter-independent program. At an intercall boundary its only persistent observer memory is its label $s\in S$. Its next command distribution is a fixed function of $s$. A physical port name, elapsed time, subprogram index, random seed or former report is not an extra free input. During a call its current command, current report and fresh coins are available. Whatever persists to another call must be encoded in $S$. A fixed decoder $r:S\to H$ reads only the committed label.

In the \emph{event-stopped} interface a post-report kernel emits $(s',e)\in S\times\{0,1\}$. An event $e=1$ ends the experiment and retains $s'$ for the decoder; no subsequent command is executed. This public commitment event is not a stored bit used in continuing computation. The same label may be used live on one path and terminal on another. In the \emph{state-stopped} interface a fixed map $h_*:S\to\{0,1\}$ determines halt status. The initial label is live, and type-one labels are globally absorbing terminal labels with no executed command rows. Entry into a terminal label is allowed only after a productive report. Thus live and terminal types are disjoint in this interface, and only in this interface by definition.

Write $R_i^\Theta(N,M)$, $i\in\{\cp,\ev,\st\}$, for the infimum of worst-case expected squared loss over programs with robust almost-sure commitment and $\sup_\theta\E_\theta C_\tau\le N$. An empty legal class has value $+\infty$. A program chosen once at random is not a free persistent seed. The ideal measurable class may have fixed real constants; effective implementation is stated separately in Section~\ref{sec:composition}.

\subsection*{A physically acquired comparison law}
Execute a declared prepared comparison experiment and obtain its joint report--audit probability $K_0(dx,dz)$ on $X\times H$. Its report marginal is $\nu$. The law is genuinely preparable and its specification has its own resource account. It need not be a member of $\Theta$. In particular no finite set of calibration data is silently promoted to a uniform certificate. Set
\begin{equation}\label{eq:derived}
 f(x)=\int zK_0(dz\mid x),\quad
 w(x)=\int\|z-f(x)\|^2K_0(dz\mid x),\quad
 d_r(x)=w(x)+\|f(x)-r\|^2.
\end{equation}
Disintegration gives Borel versions with $f\in H$; null sets may be completed arbitrarily. The profile is
\begin{equation}\label{eq:profile}
 \Q_m(b)=\min_{\substack{C\subset H,\ 1\le|C|\le m\\
                  0\le h\le b,\ \int h\,d\nu=1}}
        \int [w+\dist(f,C)^2]h\,d\nu,
 \qquad b\ge1,\ m\ge1.
\end{equation}
We use $+\infty$ for an infeasible cap or width. The quantities $f,w$ are task coordinates of acquired predictive information, not a complete substitute for the predictive quotient. Their derivation and the quotient construction in Appendix~\ref{sec:quotient} precede any use of local dimension.

\subsection*{Lower certificates on actual histories}\label{sec:lower-cert}
Fix \emph{one} actual model $\theta_*\in\Theta$, independent of the competing program. At each actual reached productive pre-call history $h$ and legal command $a$, write $P_{h,a}(dy,dz)$ for its conditional full current report--audit law. Let $P^0_{h,a}$ be a measurable comparison row on the same report space, with a Borel mark $x=\pi_{h,a}(y)$ of law $\nu$ and
\begin{equation}\label{eq:full-moments}
 \E^0[Z\mid y]=f(\pi_{h,a}y),\qquad
 \E^0[\|Z\|^2\mid y]=w(\pi_{h,a}y)+\|f(\pi_{h,a}y)\|^2.
\end{equation}
These are conditional identities given the \emph{full} report, not just its mark. Versions and inequalities are required on actual reachable histories, not only almost everywhere under an unrelated comparison path law.

Fix $0<l_-\le1\le u_-<\infty$. For all Borel $q:Y\to[0,1]$ and all finite-image Borel $r:Y\to H$, suppose predictable nonnegative defects $a_-(h,a),e_-(h,a)$ obey
\begin{align}
 \int q\,d(P_{h,a})_Y
 &\le u_-\int q\,d(P^0_{h,a})_Y+a_-(h,a),\label{eq:lower-mass}\\
 \int q(y)\|z-r(y)\|^2P_{h,a}(dy,dz)
 &\ge l_-\int q(y)\|z-r(y)\|^2P^0_{h,a}(dy,dz)-e_-(h,a).
 \label{eq:lower-task}
\end{align}
Independent current randomization can be adjoined to both rows: integrating inequalities for each fixed random value proves the same inequalities with the same defects. All command-conditioned lower inequalities hold before averaging the chosen command.

Let $A_\tau$ count productive calls. At this same actual model require, for every legal competing procedure,
\begin{equation}\label{eq:effort}
 \E_* C_\tau\ge\kappa_-\E_* A_\tau,\qquad\kappa_->0.
\end{equation}
At budget $N$ suppose, uniformly over all robustly feasible competitors,
\begin{equation}\label{eq:occupation}
 \E_*\sum_{j\le\tau:\ j\ {\rm productive}}a_-(H_{j-1},a_j)\le A_N,
 \qquad
 \E_*\sum_{j\le\tau:\ j\ {\rm productive}}e_-(H_{j-1},a_j)\le E_N.
\end{equation}
These are actual occupation certificates, not properties merely of the proposed upper program. Uniform row defects give $A_N\le(N/\kappa_-)\bar a_-$ and $E_N\le(N/\kappa_-)\bar e_-$. Lemma~\ref{lem:drift} supplies a nonuniform alternative.

\subsection*{One common upper construction}\label{sec:upper-cert}
Independently, let $\mathcal G$ be a parameter-independent program which, from each of its start histories, reaches one productive report. Its conditional expected attempt cost is at most $\kappa_+>0$ on every surviving history in every $\theta\in\Theta$. It uses $B$ live roles and returns to its start label on rejection; no reset of the physical state is assumed. It forms a virtual mark $x$, a current choice $J$ and a permitted export $T_Jr$ of its task readout. For every finite dictionary $C$ assume
\begin{equation}\label{eq:exports}
 |\{T_jr:j\in J,\ r\in C\}|\le D_0|C|,
 \qquad D_0\in\mathbb N.
\end{equation}
Every continuing phase and coin needed for this construction is included in the $B$ live roles. For fixed $0<l_+\le1\le u_+<\infty$ and $a_+,e_+\ge0$, assume for all Borel $q:X\to[0,1]$ and finite-image Borel $r:X\to H$, conditionally on every such history,
\begin{align}
 \E[q(x)]&\ge l_+\int q\,d\nu-a_+,\label{eq:upper-mass}\\
 \E[q(x)\|Z-T_Jr(x)\|^2]
 &\le u_+\int q(x)d_{r(x)}(x)\nu(dx)+e_+.
 \label{eq:upper-task}
\end{align}
The expectation here includes the within-attempt choice $J$. Unlike the lower, the upper is not required to hold for each choice separately. The choices retained after this call are nevertheless charged through the actual export or live role. These hypotheses express one common causal construction, not modelwise optimizers.

\section{Reversible response groups and common controlled memory}\label{sec:r72-reversible}
The fixed-phase results concern one transition matrix. Several commands require one family of matrices on the same retained labels, and checking every command in isolation is insufficient. We give an intrinsic finite-realization criterion for responses whose minimal linear dynamics is reversible and compact, including any finite family of unitary commands on a prepared finite-dimensional quantum system. We then obtain a matching arithmetic memory law for a pair of noncommuting involutions. This is a theorem about exact controlled response realization, not an assertion that a simulator's private label is an observed sufficient statistic for arbitrary feedback optimization.

Positive realization through invariant polyhedral sets is classical \cite{Heller65,MonrasWinter14,ThonJaeger15}. The argument below uses a section of the realizing simplex before passing to the observable quotient; auxiliary hidden-state responses need not belong to the minimal response space. Compact reversibility then turns containment into equality. The resulting finite-group obstruction concerns all command words, rather than a single spectral restriction.

\subsection{The minimal response space}
Let $A$ and $C$ be nonempty finite alphabets of commands and terminal outcomes. A supplied response is a function $b:A^*\to\Delta(C)$. Its residual after a word $w$ is $b^w(z)=b(wz)$, including the empty word $\varnothing$. Put
\[
 V_b=\operatorname{span}_{\R}\{b^w:w\in A^*\},
 \qquad (L_af)(z)=f(az).
\]
Assume $V_b$ is finite dimensional. It is invariant under each $L_a$, since $L_ab^w=b^{wa}$. Call the response \emph{compactly reversible} if the restrictions $L_a|_{V_b}$ are invertible and the closure $G_b$ of their generated group is a compact subgroup of $\operatorname{GL}(V_b)$. This concerns the observable quotient; it does not require the full physical command group to be finite or fully observable.

An exact stochastic realization with $M$ live labels has row-stochastic matrices $Q_a$, one row-stochastic terminal decoder matrix $D$, and the prescribed initial row $e_1^{\mathsf T}$, such that
\begin{equation}\label{eq:r72-response}
 b(a_1\cdots a_n)=e_1^{\mathsf T}Q_{a_1}\cdots Q_{a_n}D.
\end{equation}
Randomized terminal readout is allowed even when all $Q_a$ are deterministic. The target input word is not encoded into an uncounted clock or persistent seed.

\begin{theorem}[Finite realization of compactly reversible responses]\label{thm:r72-reversible}
For a compactly reversible response, the following are equivalent:

\textup{(i)} An exact finite stochastic realization exists.

\textup{(ii)} The residual set $\{b^w:w\in A^*\}$ is finite.

\textup{(iii)} The observable group $G_b$ is finite.

\textup{(iv)} An exact realization with deterministic transitions and a possibly randomized decoder exists.

The least number of deterministic-transition labels is the residual cardinality. This is not in general the least stochastic label count. If an $M$-label stochastic realization exists, then
\[
 |G_b|\le(2^M-1)!.
\]
This last bound is only a finite combinatorial certificate, not a sharp memory law.
\end{theorem}
\begin{proof}
Suppose \eqref{eq:r72-response} holds, and write $r_w=e_1^{\mathsf T}Q_w$. Let $R=\operatorname{span}_{\R}\{r_w\}\subseteq\R^M$. The assignment $\pi(r_w)=b^w$ extends to a well-defined surjective linear map $\pi:R\to V_b$: a relation $\sum_wc_wr_w=0$ implies, after right multiplication by $Q_zD$, the identical relation among $b^w(z)$ for every $z$. Moreover
\[
 \pi(rQ_a)=L_a\pi(r)\qquad(r\in R).
\]
The simplex section $K=R\cap\Delta_M$ is a compact polytope, invariant under right multiplication by every $Q_a$. Its image $P=\pi(K)$ is a compact polytope, contains every residual, and spans $V_b$. It satisfies $L_aP\subseteq P$. It is important that we intersect the simplex with $R$ first: projecting arbitrary hidden-state responses directly into $V_b$ would not be justified.

For an element $T$ of a compact matrix group, positive powers have a subsequence tending to the identity with exponents tending to infinity. One elementary proof takes a convergent subsequence of powers with successively growing index gaps; quotients of two close powers approach the identity. A finite-order element is covered by its period multiples. Thus positive powers also approach $T^{-1}$. Applying this to $L_a$ and using closedness of $P$ gives $L_a^{-1}P\subseteq P$, and hence $L_aP=P$. All elements of $G_b$ preserve $P$. They permute its finitely many vertices, and this permutation action is faithful because those vertices span $V_b$. Consequently $G_b$ is finite. This proves (i)$\Rightarrow$(iii).

For the displayed bound, a vertex of $K$ is determined by its positive support. If two distinct vertices had the same support, their difference could be added and subtracted in sufficiently small multiples while preserving $R$, unit mass and nonnegativity, contradicting extremality. Thus $K$ has at most $2^M-1$ vertices. A linear image of a polytope has no more vertices than the original polytope, so the faithful vertex action proves the bound.

If $G_b$ is finite, the residual orbit is finite. Conversely, if the residual set is finite, each invertible $L_a$ maps it injectively into itself, hence permutes it. Since it spans $V_b$, the generated group and its closure are finite. This proves (ii)$\Leftrightarrow$(iii). Under (ii), use the residuals as labels, send $b^w$ to $b^{wa}$ on input $a$, and use decoder $b^w(\varnothing)$. This realizes $b$ with deterministic transitions. If any such deterministic-transition machine sends $u$ and $v$ to the same label, every suffix yields the same decoder law, so $b^u=b^v$. This proves its minimal count. Finally (iv)$\Rightarrow$(i) is immediate. No assertion that stochastic transitions can always be determinized without increasing size was used.
\end{proof}

\begin{corollary}[Prepared unitary commands]\label{cor:r72-unitary}
Let $U_a$ be finitely many unitary matrices on $\mathbb C^d$, let $\rho$ be a prepared density matrix, and let $\{E_r:r\in C\}$ be a terminal positive-operator-valued measurement. For
\[
 b(w,r)=\tr(E_r\Phi_w(\rho)),\qquad \Phi_a(X)=U_aXU_a^*,
\]
exact finite stochastic realization for all command words is possible if and only if its observable response group is finite. The group is computed after quotienting only those directions invisible to every future commanded terminal test. Its dimension is at most $d^2$.
\end{corollary}
\begin{proof}
Let $W=\operatorname{span}_{\R}\{\Phi_w(\rho)\}$ in the Hermitian matrices. Each conjugation is a Hilbert--Schmidt isometry. Its forward invariant subspace $W$ is also backward invariant, by finite-dimensional invertibility. Let
\[
 N=\{X\in W:\tr(E_r\Phi_zX)=0\text{ for all }z,r\}.
\]
It is invariant under every $\Phi_a$ and hence under its inverse. The orthogonal complement of $N$ in $W$ is invariant as well. Therefore the induced transformations on $W/N$ are orthogonal for the quotient inner product. The map $X\mapsto(z,r\mapsto\tr(E_r\Phi_zX))$ identifies $W/N$ with $V_b$ and intertwines them with $L_a$. Their generated group has compact closure in the orthogonal group. Apply Theorem~\ref{thm:r72-reversible}. A constant measurement can make this induced group trivial even when the physical unitary group is infinite, so finiteness of the full physical group is not a necessary condition.
\end{proof}

\subsection{A matching memory law for two noncommuting commands}
Fix $x\in\R$ with $\sin(2\pi x)\ne0$. On a protected qubit take
\[
 U_0=\begin{pmatrix}1&0\\0&-1\end{pmatrix},\qquad
 U_1=\begin{pmatrix}\cos(\pi x)&\sin(\pi x)\\
                     \sin(\pi x)&-\cos(\pi x)\end{pmatrix},
 \quad \rho=E=|0\rangle\langle0|.
\]
Both commands are involutions. For known $0\le\alpha\le1$, their channels are
\[
 \Phi_{a,\alpha}(X)=\alpha U_aXU_a^*+(1-\alpha)\tr(X)I/2.
\]
Let $M_x^{\rm ctrl}(\alpha)$ be the least stochastic live memory of the whole binary response $w\mapsto\tr(E\Phi_{w,\alpha}\rho)$, using the same matrices and decoder for every word. In the real Bloch plane, with coordinates $(z,x)$, these commands are $\alpha S_0$ and $\alpha S_1$, where $S_0$ reflects the angle $t$ to $-t$ and $S_1$ reflects it to $2\pi x-t$. Their product is $\alpha^2R_{2\pi x}$.

\begin{theorem}[All-word arithmetic law for an involutive command pair]\label{thm:r72-control-envelope}
For $0<\epsilon<1$ and $\alpha=1-\epsilon$, the arithmetic envelope of Section~\ref{sec:r71-arithmetic} gives
\begin{equation}\label{eq:r72-control-envelope}
 \tfrac14\mathfrak A_x(\epsilon)
 \le M_x^{\rm ctrl}(1-\epsilon)
 \le(1+4\pi^2)\mathfrak A_x(\epsilon).
\end{equation}
The upper is one explicit finite stochastic machine valid for every mixed command word, not separately chosen unary machines. The lower ranges over arbitrary common stochastic realizations.

For irrational $x$, the lower and upper logarithmic growth exponents as $\epsilon\downarrow0$ are respectively
\[
 \frac1{\mu(x)+1},\qquad\frac{\mu(x)-1}{2\mu(x)-1},
\]
with values $0,1/2$ at $\mu=\infty$. A uniform cubic-root law holds precisely for badly approximable $x$; otherwise the normalized count has lower limit zero and upper limit infinity.

At $\alpha=1$, the exact count is the reduced denominator $q$ if $x$ is rational, and is infinite otherwise. At $\alpha=0$, it is two. Every strictly damped pair has a finite exact common realization. These statements concern controlled response memory, not a general multi-command boundary-optimal policy theorem.
\end{theorem}
\begin{proof}
Any common realization restricted to the words $(01)^n$ realizes
\[
 \tfrac12[1+\alpha^{2n}\cos(2\pi nx)]
\]
with transition matrix $Q_0Q_1$. Its size is therefore at least
$m_x(1-\alpha^2)\ge\mathfrak A_x(\epsilon(2-\epsilon))/2$
by Theorem~\ref{thm:r71-envelope}. For $1\le c\le2$, the defining maximum of the envelope obeys $\mathfrak A_x(c\epsilon)\ge\mathfrak A_x(\epsilon)/c$ whenever both arguments are in its domain. This follows for each denominator before taking its infimum. Taking $c=2-\epsilon$ proves the lower in \eqref{eq:r72-control-envelope}.

For the upper, the rational multiple-polygon construction in the proof of Theorem~\ref{thm:r71-envelope} gives a regular $l$-gon $P$ on the unit circle, with vertex $v_0=(1,0)$, such that
\[
 l\le(1+4\pi^2)\mathfrak A_x(\epsilon),\qquad
 \alpha R_{2\pi x}v_j\in P\quad\text{for every vertex }v_j.
\]
The first reflection $S_0$ permutes this regular polygon and zero belongs to it. Consequently $\alpha S_0P\subseteq P$, while $S_1=R_{2\pi x}S_0$ gives $\alpha S_1P\subseteq P$. Choose barycentric rows $Q_a(j,k)$ of $\alpha S_av_j$ in the same vertices. Let $d_j=(1+(v_j)_1)/2$. Starting at vertex zero, induction on the actual command word shows that the expected retained vertex equals the protected Bloch vector. Pairing with the first coordinate proves every response identity. The start is one of the counted vertices; there is no extra initialization mixture.

The two constants in \eqref{eq:r72-control-envelope} transfer both exponents and both uniform-law statements from Theorems~\ref{thm:r71-exponents} and \ref{thm:r71-cubic-rigidity}. At zero damping, a rational phase of reduced denominator $q\ge3$ has a common deterministic $q$-gon realization. The word-pair transition has a primitive $q$th-root mode and needs at least $q$ states by the unit-circle spectral obstruction. For irrational $x$ that same obstruction rules out every finite size. At $\alpha=0$ the response is one at the empty word and $1/2$ after every nonempty word, realized by a first label followed by an absorbing fair label; one label cannot change its decoder. This proves all endpoints.
\end{proof}

A fully algebraic example uses $U_1=5^{-1/2}\left(\begin{smallmatrix}2&1\\1&-2\end{smallmatrix}\right)$. Then $e^{2\pi ix}=(3+4i)/5$ is not a root of unity, since its sum with its conjugate is the rational noninteger $6/5$. Each undamped command in isolation has period at most two; their common all-word response has no finite stochastic realization. The arithmetic exponent of this particular angle is not determined here. The example isolates the failure of commandwise testing without inferring a Diophantine class from finitely many coefficients.

If an independent herald with fixed probability $p$ is placed before each command operation, and only failures apply the supplied command, every common realization above preserves the terminal response under every external command schedule and the first-herald cost $1/p$. The commands are actual inputs to that same retained-state transducer. This does not make its private label observable to the caller, allow uncharged midstream quantum measurements, or prove optimality for a controller free to choose an easier command sequence. The unary optimal squared-task interpretation below continues to use its separately proved physical embedding.

\section{A permutation normal form with an exact memory charge}
\label{sec:r73-normal-form}

Theorem~\ref{thm:r72-reversible} identifies when finite memory exists. We now
determine the form of every possible finite realization, retaining the cost of
its prescribed pure initial label. The distinction between a minimal linear
representation and a positive realization remains essential. In particular,
the permutation representation carried by the observable space does not by
itself provide a decoder with values in the probability simplex.

Let $b$ be a compactly reversible response admitting a finite realization.
Write $G$ for its finite observable group, $d=\dim V_b$, and $N$ for its residual
cardinality. We use the chronological convention $g(u)g(v)=g(uv)$: the right
action on $V_b$ is $f\cdot g(w)=L_wf$, where $L_wf(z)=f(wz)$. Thus $G$ in this
notation is the opposite of the group of operators with the usual convention
for composition; inversion identifies the two groups. No group cardinality,
faithfulness, or representation-theoretic statement depends on this convention.
Every element has a representative in $A^*$, including its inverse, since a
finite group generated by the command elements is already generated by their
positive words. The value $b(g)=b(w)$ is consequently well defined.

A \emph{probability factorization on a finite $G$-set} consists of a finite set
$X$, row permutation matrices $P_g$ with $P_gP_h=P_{gh}$, a probability row
$\lambda$ on $X$, and a row-stochastic matrix $D$ with columns indexed by $C$,
such that
\begin{equation}\label{eq:r73-probability-factorization}
 b(g)=\lambda P_gD\qquad(g\in G).
\end{equation}
All rows of $D$ are probability laws, including rows associated with states
that are not distinguished by the immediate terminal measurement. Unused
orbits may be deleted. Define
\[
 \kappa_G(b)=\min\{|X|:\text{\eqref{eq:r73-probability-factorization} holds}\}.
\]
The residual realization shows that this minimum is finite and at most $N$.
It permits a mixed $\lambda$ without charging for its preparation; the
stochastic memory $m(b)$ continues to charge every live label, beginning at
the specified point mass as in \eqref{eq:r72-response}.

\begin{theorem}[Reversible stochastic normal form]\label{thm:r73-normal-form}
Every $M$-label stochastic realization of $b$ gives a faithful finite $G$-set
$X$ and a factorization \eqref{eq:r73-probability-factorization} satisfying
\begin{equation}\label{eq:r73-initialization-charge}
 |X|+\mathbf 1_{\{\lambda\text{ is not a point mass}\}}\le M.
\end{equation}
Conversely, every such factorization has a realization with the displayed
number of labels. Its transitions on $X$ are permutations; if $\lambda$ is
mixed, its only additional label initializes the distribution on the first
command. Thus arbitrary continuing stochastic dynamics confers no advantage
over one initialization mixture followed by a permutation action.
\end{theorem}

\begin{proof}
Fix a realization $(e_1^{\mathsf T},\{Q_a\},D_0)$, and retain the reachable
row space $R$ and the surjection $\pi:R\to V_b$ from the proof of
Theorem~\ref{thm:r72-reversible}. Write
\[
 \mathcal L=\{w\in A^*:g(w)=1\}.
\]
These are identity words in the observable group; their matrices need not be
identity matrices on hidden states. Choose finitely many words
$\ell_1,\ldots,\ell_s\in\mathcal L$ whose positive matrix entries cover all
directed edges occurring in any $Q_\ell$, $\ell\in\mathcal L$. There are at
most $M^2$ distinct edges, so $s\le M^2$ is possible. If no nontrivial choice
is needed, the empty word suffices. Set
\[
 T=\frac{I+Q_{\ell_1}+\cdots+Q_{\ell_s}}{s+1}.
\]
Its closed communicating classes $C_1,\ldots,C_k$ are closed for every
$Q_\ell$ with $\ell\in\mathcal L$. The positive diagonal makes the
restriction of $T$ to each such class aperiodic. The elementary convergence
argument in Lemma~\ref{lem:r73-lazy-convergence} gives a stochastic idempotent
\[
 E=\lim_{n\to\infty}T^n,
 \qquad
 \Delta_ME=\operatorname{conv}\{\mu_1,\ldots,\mu_k\},
\]
where $\mu_i$ is the unique stationary probability row of $T$ on $C_i$.
Their supports are disjoint, so this image is a simplex with exactly $k$
vertices. Every identity-word matrix preserves each $C_i$; hence
\begin{equation}\label{eq:r73-sandwich}
 EQ_\ell E=E\qquad(\ell\in\mathcal L).
\end{equation}
This is a sandwiched equality. We do not assert $Q_\ell E=E$, $EQ_\ell=E$,
or common stationarity of the $\mu_i$ under all identity words.

Each $T^n$ is a convex combination of identity-word matrices. Consequently
$RE\subseteq R$ and $\pi(rE)=\pi(r)$ for $r\in R$. For a word $u$, put
$F_u=EQ_uE$, considered as a stochastic map of $\Delta_ME$ to itself.
Choose a word $\bar u$ representing $g(u)^{-1}$. Replacing the middle $E$
by the convex combinations $T^n$ and using
$u\ell\bar u\in\mathcal L$ gives
\[
 F_uF_{\bar u}=E=F_{\bar u}F_u.
\]
Thus $F_u$ is an affine automorphism of that simplex. Such an automorphism
permutes its vertices, because both it and its inverse preserve the simplex
and therefore preserve extremality. If $g(u)=g(v)$, the same inverse word
shows $F_uF_{\bar u}=F_vF_{\bar u}=E$, so $F_u=F_v$. Next,
$EQ_uT^nQ_vE$ is a convex combination of $F_{u\ell v}$, all equal to
$F_{uv}$. Passing to the limit proves $F_uF_v=F_{uv}$. We have obtained a
permutation action $P_g$ on the vertices $X=\{\mu_1,\ldots,\mu_k\}$.

This action is faithful. If $F_u$ is the identity on $\Delta_ME$, then
$rEQ_uE=rE$ for all rows $r$. For $r\in R$, application of $\pi$ shows
$L_u\pi(r)=\pi(r)$. Surjectivity gives $L_u=I$ on $V_b$, and so $g(u)=1$.
Write $e_1^{\mathsf T}E=\sum_i\lambda_i\mu_i$ and set
$D(i,\cdot)=\mu_iD_0$. Intertwining on $R$ gives, for every word $w$,
\[
 \lambda P_{g(w)}D=e_1^{\mathsf T}EQ_wED_0=b(w),
\]
including the empty word. If the original initial label lies in $C_i$, its
row under $E$ is exactly $\mu_i$. A mixed $\lambda$ therefore requires that
the initial label lie outside the union of the closed classes. The $k$
nonempty disjoint classes then use at most $M-1$ labels. This proves
\eqref{eq:r73-initialization-charge}; a pure $\lambda$ only requires $k\le M$.

For the converse, when $\lambda$ is a point mass its support point is the
initial label in $X$. If it is mixed, adjoin a label $\star$ with decoder
$D(\star,\cdot)=\lambda D$. On command $a$, its transition row into $X$ is
$\lambda P_{g(a)}$, and no transition returns to $\star$. The rows on $X$
are those of $P_{g(a)}$. Induction gives \eqref{eq:r73-probability-factorization}
for every nonempty word, and the decoder at $\star$ gives the empty word.
The first input is applied during initialization, so it is not lost or
encoded in an uncounted preparatory clock.
\end{proof}

The general structure of nonnegative idempotents and bounded groups of
nonnegative matrices is classical; see Brown~\cite{Brown64},
Schwarz~\cite{Schwarz64}, and Flor~\cite{Flor69}. The proof above supplies the
additional steps required for an arbitrary semigroup realization: the
identity-word projection is constructed before passing to the observable
quotient, the sandwiched action is proved to factor through that quotient,
and the prescribed pure start is charged. It does not assume that the
original stochastic matrices themselves form a group.

\begin{theorem}[Exact memory formula]\label{thm:r73-memory-formula}
With randomized terminal readout allowed and the initial label prescribed
pure, the minimum stochastic memory is
\begin{equation}\label{eq:r73-memory-formula}
 m(b)=\min_{(X,\lambda,D)}
 \left(|X|+\mathbf1_{\{\lambda\text{ is not a point mass}\}}\right)
 =\min\{N,1+\kappa_G(b)\}.
\end{equation}
The minimum ranges over precisely the probability factorizations
\eqref{eq:r73-probability-factorization}, not arbitrary linear factorizations.
\end{theorem}
\begin{proof}
The first equality is both directions of Theorem~\ref{thm:r73-normal-form}.
A point-mass factorization gives a deterministic-transition realization on
$X$, so $|X|\ge N$ by Theorem~\ref{thm:r72-reversible}. A mixed factorization
costs at least $1+\kappa_G(b)$. Conversely, the residual machine has $N$
labels, and a factorization attaining $\kappa_G(b)$ costs at most
$1+\kappa_G(b)$. These inequalities give the second equality, also when a
minimum factorization happens to have a pure start.
\end{proof}

\begin{corollary}[Sharp group bounds]\label{cor:r73-group-bound}
An $M$-label realization of a compactly reversible response satisfies
\[
 |G|\le M!,\qquad
 M<N\ \Longrightarrow\ |G|\le(M-1)!.
\]
Both bounds are attained by families below. In particular, a stochastic
memory saving over the residual count requires a faithful permutation core
on at most $M-1$ points.
\end{corollary}
\begin{proof}
The core action in Theorem~\ref{thm:r73-normal-form} embeds $G$ into $S_k$.
If $M<N$, its start cannot be pure, since that would produce a deterministic
machine with at most $M$ labels. Thus $k\le M-1$ in this case.
\end{proof}

Let $\nu_G(V_b)$ be the least cardinality of a finite $G$-set whose real
permutation module contains $V_b$, and let $\mu(G)$ be the least cardinality
of a faithful finite $G$-set. All references to $V_b$ here use the right-action
convention already fixed. A probability factorization gives $V_b$ as a
subquotient of $\mathbb R^X$: take the cyclic span of $\lambda$, then its
observable quotient. Averaging an inner product over finite $G$ splits every
invariant subspace orthogonally, so this subquotient embeds in
$\mathbb R^X$. We obtain
\begin{equation}\label{eq:r73-module-bounds}
 d\le\nu_G(V_b)\le\kappa_G(b),\qquad
 m(b)\ge\max\{d,\mu(G),\min(N,1+\nu_G(V_b))\}.
\end{equation}
These are lower bounds, not an identification of a linear module with a
bounded decoder. Proposition~\ref{prop:r73-decoder-obstruction} makes the
distinction strict.

\begin{corollary}[Rank rigidity]\label{cor:r73-rank-rigidity}
For a finite compactly reversible response,
\[
 m(b)=d\quad\Longleftrightarrow\quad N=d.
\]
In that case every realization with $d$ labels has permutation transition
matrices. If $N>d$, then $m(b)\ge d+1$.
\end{corollary}
\begin{proof}
Every $M$-label realization has $d\le M$. If $M=d$, the core in
Theorem~\ref{thm:r73-normal-form} has $k\ge d$ by its linear rank, hence
$k=d$ and its start is pure. Thus $N\le d$, whereas $N\ge d$ because the
residuals span $V_b$. Conversely $N=d$ gives a deterministic machine with
$d$ labels. For the assertion about the original matrices, a $d$-label
realization has $R=\mathbb R^d$ and the observable map $\pi$ is an
isomorphism. An inverse word for $g(a)$ therefore gives a stochastic inverse
of $Q_a$ on all of $\mathbb R^d$. A stochastic map with a stochastic inverse
is an affine automorphism of the full simplex and hence a permutation.
\end{proof}

\subsection{Two commands and a factorial deterministic gap}
Let $k\ge3$, let $P_g$ be the natural row permutation action of $S_k$, choose
a nonconstant $\lambda\in\Delta_k$, and choose a nonconstant column
$t\in[0,1]^k$. Two commands are the $k$-cycle $(1\,2\,\cdots\,k)$ and the
transposition $(1\,2)$; these generate $S_k$. Supply the binary response
\begin{equation}\label{eq:r73-symmetric-response}
 b(g)=\bigl(\lambda P_gt,\,1-\lambda P_gt\bigr).
\end{equation}
If the distinct coordinate values of $\lambda$ have multiplicities
$n_1,\ldots,n_j$, put $N_\lambda=k!/(n_1!\cdots n_j!)$.

\begin{proposition}[Exact symmetric-group family]\label{prop:r73-symmetric-family}
The response \eqref{eq:r73-symmetric-response} has observable group $S_k$,
dimension $k$, residual cardinality $N_\lambda$, and
\[
 \kappa_{S_k}(b)=k,\qquad
 m(b)=\min\{N_\lambda,k+1\}.
\]
When the coordinates of $\lambda$ are pairwise distinct, the least stochastic
memory is $k+1$ and the least deterministic-transition memory is $k!$.
For a point-mass $\lambda$, both are $k$. If either $\lambda$ or $t$ is
constant, the response instead has one residual and one-label memory.
\end{proposition}
\begin{proof}
For any nonconstant vector, differences between it and its transposes yield
a nonzero multiple of some $e_i-e_j$. Conjugating by all permutations shows
that the span of its orbit differences is
$\mathbf1^\perp$. The orbit of $\lambda$ therefore spans $\mathbb R^k$,
because every row has mass one. The orbit of $t$ also spans
$\mathbb R^k$: its differences span $\mathbf1^\perp$, and its coordinate
sum is positive, since a nonconstant nonnegative vector is nonzero.
Consequently the map
\[
 r\longmapsto\bigl(h\longmapsto(rP_ht,\,r\mathbf1-rP_ht)\bigr)
\]
is injective on $\mathbb R^k$. It identifies the reachable row space with
$V_b$, and identifies the observable action with the faithful natural
permutation action. Its residuals are precisely the distinct permutations
of $\lambda$, giving $d=k$ and $N=N_\lambda$. The displayed factorization
and the rank lower bound give $\kappa=k$. The memory formula proves the
remaining assertion; it also shows $N_\lambda\ge k$. The two constant
cases are immediate from \eqref{eq:r73-symmetric-response}.
\end{proof}

This is an exact prepared quantum experiment in every dimension $k$: take
$\rho=\operatorname{diag}(\lambda)$, permutation unitaries for the commands,
and terminal effect $E=\operatorname{diag}(t)$. The construction already
lies in the diagonal, classical subalgebra. It demonstrates a stochastic
versus deterministic memory separation, and does not assert a quantum
advantage for this family. It also proves sharpness of
Corollary~\ref{cor:r73-group-bound}: the point-mass family has $M=k$ and
$|G|=k!$, while pairwise distinct coordinates have $M=k+1<N=k!$ for $k\ge3$.

For a concrete rational example, take $k=3$,
$\lambda=(1/2,1/3,1/6)$, and $t=(0,1/4,1)^{\mathsf T}$.
The six first-outcome values are
\[
 \left\{\frac14,\frac38,\frac7{24},\frac{13}{24},
                 \frac{11}{24},\frac7{12}\right\}.
\]
It has $d=3$, $N=6$, $\kappa=3$, and $m=4$.

\subsection{The decoder constraint can remove the entire saving}
The representation-theoretic lower bound in
\eqref{eq:r73-module-bounds} need not attain the positive factorization
minimum, even for the same group, rank, and residual count as this example.

\begin{proposition}[A strict bounded-decoder obstruction]
\label{prop:r73-decoder-obstruction}
Let $u=(1,-1,0)$ and let $S_3$ act by its natural permutation matrices. The
binary response with first coordinate
\begin{equation}\label{eq:r73-endpoint-response}
 \beta(g)=\frac12+\frac14 uP_gu^{\mathsf T}
\end{equation}
has $d=3$, $N=6$, and $\nu_{S_3}(V_b)=3$, but
\[
 \kappa_{S_3}(b)=m(b)=6.
\]
It is realized by a prepared qubit with
\[
 U_0=\begin{pmatrix}1&0\\0&-1\end{pmatrix},\qquad
 U_1=\frac12\begin{pmatrix}1&\sqrt3\\\sqrt3&-1\end{pmatrix},
 \qquad \rho=E=|+\rangle\langle+|.
\]
Thus the least stochastic memory can differ for responses having the same
observable group, module, and residual cardinality.
\end{proposition}
\begin{proof}
The orbit of $u$ consists of its six distinct permutations and spans the
two-dimensional zero-sum subspace. Pairing with all permutations of $u$
is nondegenerate on that subspace. Including the constant outcome sum gives
dimension three and six distinct residuals. This representation is the
natural permutation module $\mathbb R^3$, and the rank bound gives
$\nu=3$.

Write $s=(1\,2)$ and $r=(1\,2\,3)$. Directly from
\eqref{eq:r73-endpoint-response},
\[
 \beta(1)=1,\quad\beta(s)=0,\quad
 \beta(r)=\beta(r^2)=\frac14,
 \quad\beta(\tau)=\frac34
 \quad\text{for the other two transpositions }\tau.
\]
Suppose a probability factorization exists on at most five points. Its
action is faithful, by the same residual argument as in
Theorem~\ref{thm:r73-normal-form}. Transitive $S_3$-sets have sizes
$1,2,3,6$, since their stabilizers have orders $6,3,2,1$; the subgroups of
order two are conjugate and the subgroup of order three is unique. Thus
our faithful action has one natural three-point orbit and, at most, one
two-point sign orbit together with fixed points. Two three-point orbits
already exceed the budget, and actions on only fixed and sign orbits
annihilate the alternating subgroup.

Let $a$ be any point of positive initial mass and write $q$ for the
first-coordinate decoder. The endpoint values force $q(a)=1$ and
$q(a\cdot s)=0$. No fixed point can carry mass. Within the natural orbit,
the point fixed by $s$ cannot carry mass, and both of the points exchanged
by $s$ cannot carry positive mass simultaneously. Thus the initial mass in
that orbit is concentrated at one moved point $a$; denote its total by
$A$. If a sign orbit is present its mass is concentrated at a single point;
write this total as $B$. We have $A+B=1$. Every even permutation has decoder
one on the mass in the sign orbit, and every odd permutation has decoder
zero there.

The images $a\cdot r$ and $a\cdot r^2$ are the other two points in the
natural orbit. One is $a\cdot s$, with decoder zero, and the other's decoder
is some $t\in[0,1]$. The two equal values $\beta(r)=\beta(r^2)=1/4$ give
$B=1/4$ and $At=0$. Hence $A=3/4$ and $t=0$. Of the two transpositions
different from $s$, one fixes $a$ and the other sends it to the third point.
Their first-outcome values would therefore be $A$ and zero, whereas both
must be $3/4$. This contradiction rules out all factorizations on at most
five points. The regular six-point residual machine gives $\kappa\le6$,
so $\kappa=6$ and Theorem~\ref{thm:r73-memory-formula} gives $m=6$.

For the physical realization, the adjoint actions of $U_0$ and $U_1$ on
the real Bloch plane are two reflections whose product is a rotation by
$2\pi/3$. The initial and terminal Bloch vector is the $x$-axis vector.
The first reflection sends it to its negative, the other two reflections
have inner product $1/2$ with it, and the two nonidentity rotations have
inner product $-1/2$. The Born probability is one half of one plus that
inner product. It has exactly the displayed six group values, so it is
\eqref{eq:r73-endpoint-response} after the specified generator assignment.
\end{proof}

The cyclic partition formula of the preceding revision is recovered as a
special case in Proposition~\ref{prop:r73-cyclic-module}. For the general
finite group, however, both the start charge and the bounded-decoder
constraint belong in the exact minimization problem. Neither can be
removed by a rank calculation.

\section{Exact synthesis from algebraic observable data}\label{sec:r73-algebraic}
The preceding normal form concerns existence over the real numbers. We now
give a terminating exact procedure on algebraic input, including a certificate
that no finite stochastic realization exists. The input is a supplied
finite-dimensional probability response, rather than a finite sample from an
unknown process. All algebraic numbers are specified by defining polynomials
and isolating data, with a common real number field and its chosen embedding.
The algorithms below concern exact label count and exact controller tables.
They make no polynomial running-time or bounded-bit sampling claim.

Finite matrix-group recognition is classical, and substantially faster
algorithms than the elementary enumeration below are known
\cite{BabaiBealsRockmore93,Serre07}. Reduction of probabilistic-automaton
minimization to the existential theory of the reals is also established
\cite{KieferWachter14}. Together with the inherited reversible criterion,
these ingredients already yield decidability and bounded-size minimization
on the input class considered here. The new normal form supplies the more
structured conclusion: a smallest controller can be returned as a finite
permutation core with a precisely counted initializer and a bounded
multi-output decoder. We give a complete terminating construction, including
finite certificates of the negative answer. Neither finite-group recognition
nor real-algebraic automaton minimization is asserted to be a new algorithm.

\subsection{Observable reduction and a finite stopping bound}
Write the supplied representation in chronological row convention as
\begin{equation}\label{eq:r73-algebraic-input}
 b(w)=\beta T_w B,\qquad T_{uv}=T_uT_v,
\end{equation}
where the matrices, row and decoder have entries in a real number field
$K$ of degree $e$ over $\mathbb Q$. The displayed representation is linear;
its matrices and intermediate coordinates need not be stochastic or
nonnegative. Probability validity means that the supplied $b(w)$ belongs to
$\Delta(C)$ for every word. Compact reversibility is imposed on its
\emph{minimal observable} response space, as in
Section~\ref{sec:r72-reversible}.

For clarity, minimal reduction itself is finite. Starting with $\beta$,
close its row span under right multiplication by each $T_a$. Starting with
the columns of $B$, close their span under left multiplication by the
$T_a$. Either iteration stabilizes once its dimension fails to increase,
and neither dimension exceeds the supplied dimension. A reachable row is
invisible precisely when it annihilates the second span. Quotienting the
reachable span by this kernel gives the minimal representation, with
dimension $d=\dim V_b$ and coefficients still in $K$. This is the familiar
reachable/observable linear reduction; it does not identify arbitrary
hidden stochastic rows with observable states. Since the response is
normalized and nonempty, $d\ge1$.

Set
\begin{equation}\label{eq:r73-group-budget}
 h=ed,\qquad
 B_h=|\operatorname{GL}_h(\mathbb F_3)|
     =\prod_{j=0}^{h-1}(3^h-3^j).
\end{equation}
This deliberately conservative bound has the advantage of a short complete
proof. Restriction of scalars may use a larger field than necessary; doing
so changes the computational bound, not the resulting exact label minimum.

\begin{lemma}[Classical finite-group bound]\label{lem:r73-minkowski}
A finite subgroup of $\operatorname{GL}_d(K)$ has order at most $B_{ed}$.
Consequently a monoid generated by invertible $d$-dimensional matrices over
$K$ is finite if and only if breadth-first enumeration of its distinct
positive-word products closes before producing $B_{ed}+1$ elements.
\end{lemma}
\begin{proof}
Restriction of scalars embeds the group faithfully in
$\operatorname{GL}_{ed}(\mathbb Q)$. For a finite subgroup $H$ of
$\operatorname{GL}_h(\mathbb Q)$, let
$\Lambda=\sum_{g\in H}g\mathbb Z^h$. There is a common denominator $q$
with $\mathbb Z^h\subseteq\Lambda\subseteq q^{-1}\mathbb Z^h$.
Thus $\Lambda$ is a free abelian group of rank $h$, and it is invariant
under every $g$ and $g^{-1}$. In a lattice basis, $H$ is a subgroup of
$\operatorname{GL}_h(\mathbb Z)$.

The kernel of reduction modulo $3$ has no nonidentity element of finite
order. Indeed, such an element would have a power $A$ of prime order
$\ell$. Write $A=I+3^s Z$, with $s\ge1$ and with an integer matrix $Z$
not zero modulo $3$. The binomial expansion of $A^\ell=I$, divided by
$3^s$, gives $\ell Z=0$ modulo $3$, so $\ell=3$. The expansion for this
prime, divided by $3^{s+1}$, instead gives
\[
 Z+3^sZ^2+3^{2s-1}Z^3=0,
\]
again contradicting $Z\not\equiv0\pmod3$. Reduction is therefore
injective on $H$. Counting independent ordered columns over $\mathbb F_3$
proves the bound in \eqref{eq:r73-group-budget}. This is the standard
Minkowski argument \cite{Serre07}; the proof is included to make the
termination certificate explicit.

A finite monoid of invertible matrices is a group: for every element,
two positive powers agree, and invertibility yields a positive power
equal to the identity and another equal to its inverse. A finite generated
monoid therefore has at most $B_{ed}$ elements. In a breadth-first search,
process each newly encountered product and multiply it by all generators.
If the queue empties, the finite set found is closed under all generators
and is the entire monoid. If the monoid is infinite, the search eventually
finds $B_{ed}+1$ distinct products. These are the two asserted stopping
events. All equality tests are exact tests in $K$.
\end{proof}

\begin{theorem}[Decidability and minimal algebraic synthesis]
\label{thm:r73-algebraic-synthesis}
For a supplied compactly reversible probability response with algebraic
linear representation \eqref{eq:r73-algebraic-input}, there is a terminating
algorithm that decides whether an exact finite stochastic realization
exists. If the answer is positive, the algorithm returns its least number
$m(b)$ of live labels and an exact realization with algebraic entries,
a prescribed initial label and a row-stochastic terminal decoder. It can
return a permutation core with at most one transient initializer, as in
Theorem~\ref{thm:r73-normal-form}.

A negative answer has a finite certificate consisting of $B_{ed}+1$
distinct observable command products. A positive answer has an exact
controller, the finite observable-group multiplication table, and finite
real-algebraic feasibility or infeasibility certificates for every smaller
candidate size. The negative certificate concerns the promised compactly
reversible input class; it is not a nonrealizability test for arbitrary
dissipative responses.
\end{theorem}
\begin{proof}
Perform the observable reduction and apply Lemma~\ref{lem:r73-minkowski}
to the minimal command matrices. If enumeration exceeds its bound, the
positive-word monoid and its compact group closure are infinite. Theorem
\ref{thm:r72-reversible} rules out every finite stochastic realization.
The list of distinct products, each supplied with its generating word,
is a finite independently checkable certificate of this conclusion.

Otherwise enumeration closes to a finite group $G$. Compute the distinct
reachable rows $\beta T_g$ in the minimal representation; their count $n$
is the residual cardinality. The deterministic residual machine proves
$m(b)\le n$. Theorem~\ref{thm:r73-memory-formula} expresses the answer as
a minimum over finite permutation actions, a probability row $\lambda$,
and a stochastic decoder $D$. Only the following finite possibilities can
improve or attain the upper $n$: a core of size $k\le n$ with a point
initial row, or a core of size $k\le n-1$ with a possibly mixed initial
row and one additional initializer.

For each such $k$, enumerate assignments of permutations to the elements
of $G$ and retain those satisfying
$P_1=I$ and $P_gP_h=P_{gh}$. This is a finite search through at most
$(k!)^{|G|}$ assignments, and it may of course be reduced using generators
and relations. For a retained action impose
\begin{equation}\label{eq:r73-bilinear-synthesis}
 \lambda_i\ge0,\quad\sum_i\lambda_i=1,\qquad
 D_{ic}\ge0,\quad\sum_cD_{ic}=1,\qquad
 \lambda P_gD=b(g)\quad(g\in G).
\end{equation}
For the point case impose $\lambda=e_i^{\mathsf T}$ and enumerate $i$.
For the other case it is harmless to permit point rows as well: that
candidate then pays a redundant initializer, and the separately enumerated
point case can only improve its cost. Equations
\eqref{eq:r73-bilinear-synthesis} have algebraic coefficients and degree
at most two in their real unknowns. Exact real-algebraic elimination
decides feasibility, and a nonempty such set has an algebraic sample point.
One may see the latter directly from a sign-invariant algebraic cell
decomposition: choose isolating algebraic sample coordinates successively
in its nonempty sections and sectors. Algebraic elimination is the
classical decision theorem used here \cite{Seidenberg54}; no numerical
optimizer is being substituted for it.

Search in increasing total label cost and stop at the first feasible
candidate. Termination follows from the deterministic upper. Theorem
\ref{thm:r73-memory-formula} proves that no stochastic competitor outside
the enumerated normal forms can be smaller. For a point initial row use
the permutation matrices directly. For a mixed initial row add the
counted initial state, give it decoder $b(\varnothing)=\lambda D$, and
on command $a$ send it to $\lambda P_{g_a}$ in the core. All later
transitions are permutations. Every returned entry is algebraic, and all
response identities follow from the group action and
\eqref{eq:r73-bilinear-synthesis}.

The finite word criterion of Theorem~\ref{thm:r73-finite-tests} supplies
an alternative bounded-size formulation using arbitrary stochastic
matrices, independently of the normal form. This is useful for checking
a returned implementation against the original linear representation.
Its existence does not give an efficient algorithm for minimizing
probabilistic automata, a problem with established hardness results
\cite{KieferWachter14}.
\end{proof}

\begin{corollary}[Algebraic prepared unitary experiments]
\label{cor:r73-unitary-algorithm}
Let the prepared density matrix, a finite family of unitary commands and
a finite terminal positive-operator-valued measurement have algebraic
real and imaginary entries. Whether their Born probabilities for all
command words admit an exact finite classical stochastic realization is
decidable. In the positive case a smallest exact classical realization
and its bounded decoder can be synthesized with algebraic entries.
For a protected Hilbert space of dimension $d_{\rm phys}$, the observable
linear dimension used in the algorithm is at most $d_{\rm phys}^2$.
\end{corollary}
\begin{proof}
Hermitian matrices have a real coordinate basis of size
$d_{\rm phys}^2$ consisting of diagonal matrix units and real and
imaginary symmetric matrix-unit pairs. In this basis conjugation by each
supplied unitary, preparation, and terminal evaluation have real algebraic
coefficients. This constructs \eqref{eq:r73-algebraic-input}. Positivity
and normalization of Born probabilities follow from the declared density
matrix and measurement, not from a finite test. The observable reduction
is finite linear algebra over a common real number field. Corollary
\ref{cor:r72-unitary} proves compact reversibility of this quotient, so
Theorem~\ref{thm:r73-algebraic-synthesis} applies.
\end{proof}

This corollary allows noncommuting commands, nonfaithful physical actions,
mixed preparations and arbitrary finite terminal measurements. In
particular, a physically infinite group with constant observable response
is correctly reduced to the one-label case before group enumeration. The
result concerns exact simulation of the supplied experiment. Choosing an
optimal feedback experiment, deciding validity of an arbitrary algebraic
series, approximate program size, exact sampling and workspace lower
bounds are different questions. The complete retained causal and boundary
theory keeps its own quantifiers and resource accounts.

\section{Finite certificates and stable exact memory}\label{sec:r73-finite-certificates}

The normal form gives exact realizations at reversible responses. We next prove
that their memory obstructions admit finite tests and persist under a specified
class of physical perturbations. Finite-word equivalence for linear machines is
classical \cite{Tzeng92}; the finite bound below is its direct difference-space
argument. The additional conclusions concern the positive separation from
\emph{all} smaller stochastic machines, its effective algebraic certificate,
and a matching realization under damping toward a common background response.

\subsection{Finite tests against a fixed memory budget}
Suppose a response $b:A^*\to\Delta(C)$ has a supplied real linear representation
\[
 b(w)=\beta T_w B,
 \qquad \beta\in\R^{1\times d},\quad T_a\in\R^{d\times d},
 \quad B\in\R^{d\times |C|}.
\]
Minimality of this representation is not required. For $M\ge1$ let
$\mathcal P_M$ be the compact product of the row-stochastic matrix simplexes
for $(Q_a)_{a\in A}$ and $D\in\R^{M\times |C|}$. Its initial row is the
prescribed pure row $e_1^{\mathsf T}$, as before. Set
\begin{equation}\label{eq:r73-finite-gap}
 L=d+M-1,\qquad
 \delta_M(b)=\min_{(Q,D)\in\mathcal P_M}
 \max_{\substack{|w|\le L\\c\in C}}
 \left|e_1^{\mathsf T}Q_wD_c-b(w,c)\right|.
\end{equation}

\begin{theorem}[Finite obstruction and effective separation]\label{thm:r73-finite-tests}
An $M$-label stochastic machine realizes $b$ on every word if and only if it
does so on the words of length at most $d+M-1$. Consequently
$\delta_M(b)>0$ if and only if no exact $M$-label realization exists.
If another supplied response $\widetilde b$ satisfies
\[
 \max_{|w|\le L,\ c\in C}|\widetilde b(w,c)-b(w,c)|\le\eta,
\]
then every $M$-label machine has error at least $\delta_M(b)-\eta$ on these
same probes. This last assertion imposes no finite-rank hypothesis on
$\widetilde b$.

If the given entries of $\beta,T_a,B$ are real algebraic numbers,
$\delta_M(b)$ is an effectively computable real algebraic number, and a
minimizing machine with real algebraic entries can be computed. When it is
positive, a positive rational number strictly below it and a real-algebraic
quantifier-elimination certificate of that strict lower bound can be computed.
No polynomial-time or practical-size bound is asserted.
\end{theorem}
\begin{proof}
For a fixed candidate machine use the block-diagonal matrices
$S_a=\operatorname{diag}(T_a,Q_a)$, the row
$u=(\beta,-e_1^{\mathsf T})$, and the output columns
$v_c=(B_c,D_c)^{\mathsf T}$. The difference on $w$ is $uS_wv_c$.
Let $R_j$ be the span of $uS_w$ for $|w|\le j$. These spaces increase;
if $R_j=R_{j+1}$, the common space is invariant under every $S_a$ and all
later spaces coincide. Since $u\ne0$, the sequence starts in dimension one
inside a space of dimension $d+M$ and must have stabilized by $j=d+M-1$.
Vanishing on the indicated finite set therefore implies vanishing on all
words. This is the same linear-equivalence principle as the inherited
finite-realization test.

The maximum in \eqref{eq:r73-finite-gap} is a continuous function on the
compact set $\mathcal P_M$, so its minimum is attained. It is zero precisely
when one machine passes every finite probe, hence, by the first assertion,
realizes every word. The perturbation claim follows pointwise from the triangle
inequality and then from the minimum defining $\delta_M(b)$.

For algebraic input, every probe value is algebraic. The constraints defining
row-stochasticity and the inequalities between the finitely many polynomial
errors and a variable $t$ are polynomial constraints over the real algebraic
numbers. Real quantifier elimination decides their feasibility and computes
the least feasible $t$ together with an algebraic witness. A singleton
semialgebraic set over the real algebraic numbers consists of an algebraic
number. If $\delta_M(b)>0$, effective isolation supplies a positive rational
$\sigma<\delta_M(b)$; infeasibility with all errors at most $\sigma$ is a
finite real-algebraic certificate. This uses the classical decision procedure
\cite{Seidenberg54}; it does not assert that this procedure was executed by
the finite diagnostics accompanying this article.
\end{proof}

For the prepared quantum response, these finite probes have an explicit
physical perturbation bound. Suppose two density matrices differ in trace
norm by at most $h_\rho$, corresponding channels differ in diamond norm by
at most $h_\Phi$, and corresponding terminal effects differ in operator norm
by at most $h_E$. All channels are trace preserving and completely positive,
and both lists of effects are valid measurements. Contractivity and telescoping
give $\|\widetilde\Phi_w(\widetilde\rho)-\Phi_w(\rho)\|_1
\le h_\rho+|w|h_\Phi$. The state difference has trace zero; subtracting
$I/2$ from a terminal effect, whose spectrum lies in $[0,1]$, gives
\begin{equation}\label{eq:r73-physical-probes}
 \max_{|w|\le L,c}|\widetilde b(w,c)-b(w,c)|
 \le h_E+\tfrac12(h_\rho+Lh_\Phi).
\end{equation}
Thus a strict smaller-memory obstruction survives a specified open
neighbourhood of physical primitives. This supplies a lower bound there;
it does not assert that an arbitrary perturbation has a realization of the
same size.

\subsection{A matching upper under a common background}
Let $b$ be compactly reversible, nonconstant, and have finite stochastic
memory $m=m(b)$. Let $G$ be its finite observable group, with the chronological
convention of the normal-form theorem, and put
\[
 \overline b=|G|^{-1}\sum_{g\in G} b(g),\qquad
 C_b=\max_{g\in G,c\in C}|b(g,c)-\overline b(c)|>0.
\]
Choose a damping factor $\alpha_a\in[0,1]$ for each command, put
$\beta_w=\prod_{j=1}^{|w|}\alpha_{w_j}$, with $\beta_\varnothing=1$, and define
\begin{equation}\label{eq:r73-common-background}
 b_{\boldsymbol\alpha}(w)
   =\beta_w b(w)+(1-\beta_w)\overline b.
\end{equation}

\begin{theorem}[Stable exact memory under background damping]\label{thm:r73-memory-plateau}
The response \eqref{eq:r73-common-background} has an exact realization with
$m$ labels for every choice of the damping factors. Put $L=d+m-2$, where
$d$ is any supplied representation dimension of $b$, and let
$\epsilon=\max_a(1-\alpha_a)$. Then
\begin{equation}\label{eq:r73-memory-plateau}
 \epsilon<\min\left\{1,\frac{\delta_{m-1}(b)}{LC_b}\right\}
 \quad\Longrightarrow\quad m(b_{\boldsymbol\alpha})=m.
\end{equation}
For algebraic input satisfying the hypotheses, a positive algebraic cutoff,
or a smaller positive rational cutoff, is effectively computable.
For a constant response every damping leaves that response unchanged and its
minimum is one label; the expression involving $m-1$ is not used.
\end{theorem}
\begin{proof}
Take a minimal realization in the normal form proved above. On its recurrent
labels each $Q_a$ is a permutation, and an extra initializer is present only
when the initial distribution on these labels is mixed. Average that
distribution over $G$, and call the resulting row $\pi$, extended by zero
on the initializer when present. It obeys
\[
 \pi Q_a=\pi,\qquad \pi D=\overline b.
\]
Indeed, the group average is invariant under each permutation. The initializer
has no incoming transition, and its first-input row is the original initial
distribution followed by the commanded permutation, so the same identities
hold on the entire machine. Set $J=\mathbf1\pi$. Then
$J^2=J$, $Q_aJ=JQ_a=J$, and
\[
 Q_{a,\boldsymbol\alpha}=\alpha_aQ_a+(1-\alpha_a)J
\]
is row-stochastic on the same labels. Multiplication gives, for every nonempty
word and also trivially at the empty word,
\[
 Q_{w,\boldsymbol\alpha}=\beta_wQ_w+(1-\beta_w)J.
\]
Reading with the same decoder proves the upper.

By minimality $\delta_{m-1}(b)>0$. For $|w|\le L$, the elementary product
bound $1-\beta_w\le |w|\epsilon$ gives an error from $b$ of at most
$LC_b\epsilon$. Theorem~\ref{thm:r73-finite-tests} rules out every
$(m-1)$-label realization when this is smaller than $\delta_{m-1}(b)$.
Machines with fewer labels can be padded by unreachable rows, so they too are
excluded. This and the upper prove equality. A nonconstant response has
$m\ge2$ and $d\ge2$, hence $L>0$. The effective assertion follows from the
finite group enumeration and algebraic synthesis theorem, together with
Theorem~\ref{thm:r73-finite-tests}. Finite maxima and averages over the
computed group preserve algebraicity.
\end{proof}

There is a physical realization of the common background for prepared unitary
commands. Let $\mathcal K$ be the compact closure of their physical conjugation
group. The compact convex hull of $\{\Phi(\rho):\Phi\in\mathcal K\}$ has a
unique element $\rho_*$ of least Hilbert--Schmidt norm: strict convexity gives
uniqueness. Every conjugation in $\mathcal K$ preserves that set and norm, so
$\rho_*$ is invariant. Its observable image belongs to the convex hull of the
finite orbit of $b$ and is invariant. Averaging any element of that orbit
over $G$ gives $\overline b$; applying this averaging projection to an
invariant element in its convex hull shows that its observable image is
exactly $\overline b$. Therefore the physical channels
\[
 \widetilde\Phi_a(X)=\alpha_aU_aXU_a^*
                   +(1-\alpha_a)\operatorname{tr}(X)\rho_*
\]
produce \eqref{eq:r73-common-background} under every command word. In general
$\rho_*$ need not be maximally mixed. The theorem concerns damping toward
this common invariant background, not every depolarizing or noisy channel.

\subsection{Periodic plateaux and an explicit rational certificate}
For the periodic mixture of Section~\ref{sec:r72-periodic}, let the distinct
nonreal frequencies have positive amplitudes $\gamma_j$, total
$\Gamma=\sum_j\gamma_j$, and let $M_*$ be its exact partition-and-anchor
minimum. Its dimension is $d=2r+1$ and $C_b\le\Gamma/2$. Hence a common
damping $\alpha=1-\epsilon$ retains exactly $M_*$ labels whenever
\begin{equation}\label{eq:r73-periodic-cutoff}
 0\le\epsilon<
 \min\left\{1,\frac{2\delta_{M_*-1}(b)}{\Gamma(2r+M_*-1)}\right\}.
\end{equation}
This is an effective finite certificate for each algebraic periodic input.
It closes the previously existential plateau cutoff for this class, without
claiming a sharp transition point or a fixed-memory risk asymptotic.

For clarity, one certificate can be calculated without a general
quantifier-elimination run. Consider the unary response
\begin{equation}\label{eq:r73-period-four}
 b_\alpha(n)=\tfrac12\{1+\alpha^n\cos(\pi n/2)\},\qquad n\ge0.
\end{equation}

\begin{proposition}[A certified four-label interval]\label{prop:r73-period-four}
For the undamped sequence $b=b_1$ and the dimension $d=3$,
$\delta_3(b)>1/36$. In particular the exact stochastic memory of
\eqref{eq:r73-period-four} is four for $89/90\le\alpha\le1$.
This interval is a certified sufficient interval, not a claimed sharp one.
\end{proposition}
\begin{proof}
Let $Q$ be any row-stochastic $3\times3$ matrix, put $t=\operatorname{tr}Q$
and $h=\det Q$, and write $y_n=e_1^{\mathsf T}Q^nD$ for any scalar decoder.
Since $Q\mathbf1=\mathbf1$, its characteristic polynomial is
\[
 z^3-tz^2+(t+h-1)z-h.
\]
Thus $y_{n+3}-ty_{n+2}+(t+h-1)y_{n+1}-hy_n=0$. We have $0\le t\le3$,
$|h|\le1$, and $|t+h-1|\le3$. The determinant bound follows directly by
bounding its absolute permutation sum by the product of the row sums.
If the first five values of $y$ are within $\eta$ of
$(1,1/2,0,1/2,1)$, substitution at $n=0,1$ yields
\[
 |t-h|\le16\eta,\qquad |2-t-h|\le16\eta,
 \qquad |t-1|,|h-1|\le16\eta.
\]
Nonnegativity and the row sums imply a further exact identity. For
$a_i=1-Q_{ii}$ it is
\begin{equation}\label{eq:r73-three-state-certificate}
 t^2-3t+3-3h
 =\tfrac12\sum_{i<j}(a_i-a_j)^2
       +3\sum_{i<j}Q_{ij}Q_{ji}\ge0.
\end{equation}
To verify it, the second characteristic coefficient is both $t+h-1$ and
$\sum_{i<j}(Q_{ii}Q_{jj}-Q_{ij}Q_{ji})$; substitute
$Q_{ii}=1-a_i$ and simplify. On the other hand, putting $u=t-1$ and $v=h-1$
gives
\[
 t^2-3t+3-3h=-2-u+u^2-3v
       \le-2+64\eta+256\eta^2.
\]
For $\eta\le1/36$ the right side is at most $-2/81<0$, a contradiction.
The compact minimum on the first five probes is therefore strictly greater
than $1/36$, and so is \eqref{eq:r73-finite-gap}, which here uses the six
probes $0\le n\le5$.

The undamped response has dimension three and a four-cycle realization with
decoder $(1,1/2,0,1/2)$. The same polynomial certificate excludes three labels,
so its minimum is four. Damping that cycle toward the uniform stationary row
realizes \eqref{eq:r73-period-four} with four labels. Its difference from the
undamped sequence on $0\le n\le5$ is at most
$\tfrac52(1-\alpha)$. For $\alpha\ge89/90$ this is at most $1/36$, strictly
below $\delta_3(b)$. The finite obstruction proves the lower and completes
the proof.
\end{proof}

The certificate uses a stochastic-matrix inequality rather than a sample
search over three-label machines. Its accompanying rational diagnostics check
the displayed polynomial identity, the constructed all-word damping
identities on explicit finite models, and exact finite recurrence consequences.
They do not replace the proofs, compute all algebraic gaps, or give an
independent human or editorial assessment.

\section{Relation to the retained exact resource laws}
\label{sec:r73-consequences}
The finite-group normal form has a cyclic specialization, but it does not
replace the independent spectral and arithmetic arguments for responses
outside the reversible class. We record the precise relation.

For the probability-valid periodic mixture of
Theorem~\ref{thm:r72-periodic-memory}, with distinct nontrivial frequency
orders $q_1,\ldots,q_r$, let
\[
 K=\min_{\mathcal P}\sum_{B\in\mathcal P}
                  \operatorname{lcm}\{q_j:j\in B\},
\]
where $\mathcal P$ ranges over partitions of the frequency indices.
The exact pure-start count is the minimum of the full period
$\operatorname{lcm}(q_1,\ldots,q_r)$ and $K+1$, with the constant and
single-frequency endpoints as stated in the retained theorem.
A cyclic permutation representation is a disjoint union of cycles; a
frequency of order $q$ occurs on a cycle precisely when $q$ divides its
length. Assigning required modes to cycles therefore gives the same
partition expression. The bounded decoders and positive-amplitude
preparation in the retained theorem attain that lower bound. This agreement
identifies a model-specific closed form inside the general factorization
problem. Arbitrary bounded decoders do not have this same closed formula,
as the strict $S_3$ obstruction shows.

The leakage theorem, Theorem~\ref{thm:r72-periodic-leakage}, adds its
separately proved inactive live label. Its lower bound uses a
decay-normalized observable core and the no-critical-chain argument;
the common-background plateau of Section~\ref{sec:r73-finite-certificates}
is a different perturbation. The exact spectral construction of
Theorem~\ref{thm:r69-spectral}, the relative-spectrum obstruction and
the arbitrary-phase envelope of Theorem~\ref{thm:r71-envelope} remain
proved with their own hypotheses. In particular a strictly damped
irrational mode need not have compactly reversible observable dynamics.

The following actual-history transfer theorem remains a general statement
about causal experiment comparison. It compares complete procedures,
charges the common paid acquisition construction, and forms its lower
witness on actual competing histories. None of its hypotheses is inferred
from a finite observable group. The complete paired supplement retains the
associated performed-law minimax and dual theory, finite-width transfer,
reset control, optimal discrimination, predictive quotients and geometric
results. The original ordered-measurement and two-stage-reset realizations
remain independently identifiable by source commit and path.

\section{The causal task-order transfer theorem}\label{sec:central}
Suppose that every terminating state-stopped competitor reaches at least $B_0\ge1$ live labels under $\theta_*$. Such a certificate must concern actual physical command compatibility, not printed action names. Define
\begin{equation}\label{eq:widths}
\begin{array}{c|cc}
 i&m_-(M)&m_+(M)\\\hline
 \cp&M&\lfloor M/D_0\rfloor\\
 \ev&M&\lfloor M/D_0\rfloor\\
 \st&M-B_0&\lfloor(M-B)/D_0\rfloor .
\end{array}
\end{equation}
The event upper also requires $M\ge B$. The trivial state-stopped lower $B_0=1$ is valid for a procedure initially waiting to make a call.

\begin{theorem}[Affine task-order transfer with a paid reference and common construction]\label{thm:transfer}
Assume all the prepared-kernel, full-report, task-order, effort, occupation, resource and common-construction hypotheses of Section~\ref{sec:data}. Fix $N,M$, put $b_-=N/\kappa_-$, $b_+=N/\kappa_+$, and suppose $b_-\ge1$ and $m_-(M)\ge1$. Set
\[
 t_N=(1-A_N)_+/u_-.
\]
If $t_N>0$, every one of the three risks satisfies
\begin{equation}\label{eq:central-lower}
 R_i^\Theta(N,M)\ge
 \bigl[l_-t_N\Q_{m_-(M)}(b_-/t_N)-E_N\bigr]_+.
\end{equation}
If $t_N=0$, the asserted lower is zero, with no evaluation of $\Q$ at infinity. An empty legal class still has risk $+\infty$.

Suppose also $m_+(M)\ge1$, $b_+\ge1$, and, for event stopping, $M\ge B$. If
\begin{equation}\label{eq:beta}
 \beta_N=\frac{l_+b_+}{1+b_+a_+}\ge1,
\end{equation}
one common Borel program with actual expected cost at most $N$ in every model attains
\begin{equation}\label{eq:central-upper}
 R_i^\Theta(N,M)\le
 \frac{u_+(1+b_+a_+)}{l_+}\Q_{m_+(M)}(\beta_N)+b_+e_+.
\end{equation}
It commits almost surely. If $\beta_N<1$, the always-accept first-attempt construction instead gives $R_i^\Theta(N,M)\le u_+\Q_{m_+(M)}(1)+e_+$ at cost at most $\kappa_+\le N$.

The lower is against every legal competing strategy and uses one actual witness. The upper is one parameter-independent strategy for all models. The two effort constants need not agree. Neither direction assumes that the comparison experiment belongs to $\Theta$ or that a normalized predictive update is uniformly contracting.
\end{theorem}

\begin{corollary}[Scale-uniform relative transfer]\label{cor:relative}
Under Theorem~\ref{thm:transfer}, suppose $A_N=E_N=a_+=e_+=0$. If $l_+b_+\ge1$, then
\begin{equation}\label{eq:relative}
 \frac{l_-}{u_-}\Q_{m_-(M)}(u_-b_-)
 \le R_i^\Theta(N,M)
 \le\frac{u_+}{l_+}\Q_{m_+(M)}(l_+b_+).
\end{equation}
The factors $l_\pm,u_\pm$ may be fixed nonunit constants, independent of budget and resolution. In particular no term proportional to $b_\pm$ times a fixed absolute calibration error is required in this relative-order class.
\end{corollary}
\begin{proof}
Substitute $t_N=1/u_-$ and $\beta_N=l_+b_+$ in the theorem.
\end{proof}

\begin{corollary}[Separate exact checkpoint and autonomous laws]\label{cor:exact}
Suppose the certificates are exact: $l_\pm=u_\pm=1$, all defects vanish, $\kappa_-=\kappa_+=\kappa$, and exports are the identity. Then, for $N\ge\kappa$,
\[
 R_\cp^\Theta(N,m)=\Q_m(N/\kappa).
\]
The constructed autonomous upper uses $B+m$ state-stopped labels, or $\max(B,m)$ event-stopped labels. If the mandatory live count is sharp, $B_0=B$, then
\begin{align}
 R_\st^\Theta(N,B+m)&=\Q_m(N/\kappa)\quad(m\ge1),\label{eq:typed-exact}\\
 R_\ev^\Theta(N,M)&=\Q_M(N/\kappa)\quad(M\ge B).\label{eq:event-exact}
\end{align}
The subtraction by $B$ belongs only to the state-stopped model. Nonidentity exports require the expanded dictionary in \eqref{eq:widths}, unless an optimal dictionary is invariant under all exports. These exact stopping distinctions are inherited from the complete companion and are reproved here to fix the resource semantics.
\end{corollary}
\begin{proof}
The same lower and upper profile coincide in each asserted interface. Identity exports give $D_0=1$. The packing and converse widths are proved in Lemma~\ref{lem:roles}. No lower on live states is used to forbid event-stopped reuse.
\end{proof}

The familiar absolute-defect theorem is also a specialization: take $l_\pm=u_\pm=1$ while retaining $A_N,E_N,a_+,e_+$. Thus relative and absolute calibration are composed in one statement. The chosen metric is the common scored task throughout; it is not changed between the converse and the construction.

\section{Proof of the central theorem}
\begin{lemma}[Profile attainment]\label{lem:attainment}
For every finite $m\ge1$ and $b\ge1$, the minimum in \eqref{eq:profile} is attained. The profile is nonincreasing in $m,b$ and left-continuous in $b>1$.
\end{lemma}
\begin{proof}
Write a dictionary as an ordered tuple in the compact set $H^m$, permitting repetitions. The feasible densities are a weak-star closed subset of the radius-$b$ ball of $L^\infty(\nu)$: positivity, the upper bound and the mass equality are each closed under pairing with $L^1(\nu)$. They are therefore weak-star compact. As the tuple varies, its squared-distance function varies uniformly on $H$, since all points belong to a bounded set. Along a convergent product net, split the objective difference into this uniform variation and the weak-star pairing against the fixed bounded measurable loss at the limit tuple. This proves continuity of the objective and attainment. Nets avoid an unwarranted sequential compactness assumption. Monotonicity follows from inclusion of feasible sets. For left continuity take an optimizer $h$ at $b>1$ and let $h_\lambda=(1-\lambda)h+\lambda$. It has unit mass and cap $b_\lambda=(1-\lambda)b+\lambda<b$. Its objective tends to the optimum as $\lambda\downarrow0$, while monotonicity bounds the optimum at $b_\lambda$ between the optimum at $b$ and this objective.
\end{proof}

\begin{lemma}[Actual-history comparison measure]\label{lem:comparison}
Fix any robustly admissible competitor of cost at most $N$, with its fixed terminal dictionary $C$. There are finite nonnegative measures $\xi_c$ on $X$, $c\in C$, such that $\xi=\sum_c\xi_c$ obeys
\begin{align}
 \xi&\le (\E_*A_\tau)\nu\le b_-\nu,\label{eq:measure-cap}\\
 1&\le u_-\xi(X)+A_N,\label{eq:measure-mass}\\
 \E_*\|Z-\widehat r\|^2
 &\ge l_-\sum_{c\in C}\int d_c\,d\xi_c-E_N.
 \label{eq:measure-risk}
\end{align}
The measure $\xi$ is not asserted to be the terminal law of any comparison execution.
\end{lemma}
\begin{proof}
Adjoin an independent uniform random variable to the current report on both sides. This is current-call randomness, not a seed retained between calls; it is sufficient to represent the competitor's randomized output and stopping decision. A finite-output kernel has such a representation by cumulative probabilities. At a fixed reached productive history $h$ and chosen command $a$, the post-report decision then gives an acceptance indicator $q_h(y)$ and an output $r_h(y)\in C$ on acceptance. Give $r_h$ an arbitrary fixed value on rejection. The notation $y$ now includes the current independent random variable. Integrating the original certificates over that variable preserves their constants. Dependence on the entire earlier history is permitted in these functions.

For each Borel $D\subset X$, set
\begin{equation}\label{eq:xi-definition}
 \xi_c(D)=\E_*\sum_{j\le\tau:\ j\ {
m productive}}
 \int\one_{\{\pi_{H_{j-1},a_j}(y)\in D\}}
 q_{H_{j-1}}(y)\one_{\{r_{H_{j-1}}(y)=c\}}
 P^0_{H_{j-1},a_j}(dy,H).
\end{equation}
The outer expectation is the actual pre-call occupation at $\theta_*$. All terms are jointly measurable and nonnegative. Tonelli defines the measures, whose total is finite because $\E_*A_\tau\le b_-$. Summing over $c$ and dropping the acceptance indicator gives \eqref{eq:measure-cap}: each comparison row has mark law $\nu$.

The analogous sum with the actual rows instead of $P^0$ has mass exactly one, because the actual procedure commits exactly once almost surely. Apply \eqref{eq:lower-mass} at each actual reached history, then sum. The occupation certificate yields \eqref{eq:measure-mass}. Apply \eqref{eq:lower-task} in the same way to the disjoint actual commitment events. Their expected losses sum to the actual risk. The comparison contribution is $d_{r_h(y)}(\pi_{h,a}y)$ by \eqref{eq:full-moments}; this identity conditions on the full current report and therefore does not overlook extra information seen by the competitor. The accumulated loss defect is at most $E_N$, giving \eqref{eq:measure-risk}. One may first truncate the number of calls. Monotone convergence for the nonnegative contributions and the finite defect bounds justify the limit.

No actual past is replaced by a comparison past. In particular the competitor need not terminate, or obey a cost bound, if hypothetically run under the comparison experiment. This distinction is why the argument applies when that experiment is not an operating model.
\end{proof}

\begin{remark}[Why the full report is required]
A constant mark can conceal a visible fair bit which also equals the future audit. A comparison using only that mark has $f=0,w=1$, but a two-value decoder reading the bit predicts it with zero loss. The hypotheses therefore condition on the full report seen by the competitor; matching a compressed mark marginal alone cannot establish the lower.
\end{remark}

\begin{lemma}[Live roles, output roles and physical compatibility]\label{lem:roles}
A program with $\ell$ live roles and $t$ output roles has a state-stopped implementation with $\ell+t$ labels and an event-stopped implementation with $\max(\ell,t)$ labels, preserving every path cost and complete stopped scored law. Conversely an event-stopped $M$-label program has at most $M$ decoder values. A state-stopped $M$-label program with at least $B_0$ reached live labels has at most $M-B_0$ terminal decoder values.

For a physical live-state certificate, consider finitely many mandatory pre-call contexts $I$ with no free context input. Suppose the commands legal at a history of context $i$ are exactly $\mathcal U_i$, and every terminating competitor visits every context with positive probability under the same witness. This positive probability may depend on the competitor; no uniform lower bound over competitors is required. Let $b_{\rm phys}$ be the least size of a cover of $I$ by subsets $J$ with $\bigcap_{i\in J}\mathcal U_i\ne\varnothing$. Then at least $b_{\rm phys}$ live labels are required. This certificate is invariant under duplicate names and bijective stateless recodings of the same physical commands.
\end{lemma}
\begin{proof}
For state stopping take disjoint copies of the live and terminal role sets and type them zero and one. For event stopping embed the live and terminal role sets separately into one alphabet of size $\max(\ell,t)$. On continuing transitions use the live embedding; on commitment emit the final event and use the terminal embedding. A command row is read only on a continuing branch, while the decoder is read only after commitment. Sharing their label therefore causes no conflict and uses no continuing mode bit. Couple the original and recoded coins. Every finite stopped prefix and every path cost agree, and almost-sure termination gives equality of the stopped scored laws. A fixed decoder has one value per retained label. In state stopping the global halt map excludes each reached live label from the output set, proving its subtraction. Event stopping has no such exclusion.

A live label reached in contexts $J$ must have one command distribution supported on all the corresponding $\mathcal U_i$. Since $J$ is finite, the union of the exceptional null command sets is null; hence their intersection supports a probability and is nonempty. The sets of contexts served by the labels cover $I$, proving the bound. Conversely a nonempty intersection contains a command and hence a Dirac command distribution. Replacing names without changing executed kernels, legal domains, costs and reports does not change these intersections. A translator which itself remembers a hidden phase is not such a recoding; its memory must be charged.
\end{proof}

\begin{proof}[Proof of Theorem~\ref{thm:transfer}]
Fix a competing program. Lemma~\ref{lem:roles} bounds its decoder width by $m_-(M)$. Let $s=\xi(X)$ in Lemma~\ref{lem:comparison}. If $t_N>0$, then \eqref{eq:measure-mass} gives $s\ge t_N$. Scale the measure to $\xi'=(t_N/s)\xi$. It has mass $t_N$ and remains bounded above by $b_-\nu$. Thus $h=(d\xi'/d\nu)/t_N$ has unit mass and cap $b_-/t_N$. Nearest-dictionary projection and nonnegativity imply
\[
 \sum_{c\in C}\int d_c\,d\xi_c
 \ge\int [w+\dist(f,C)^2]\,d\xi
 \ge t_N\Q_{m_-(M)}(b_-/t_N).
\]
Together with \eqref{eq:measure-risk} this proves the asserted lower for this competitor. Its worst-model risk dominates its risk at $\theta_*$. Taking the infimum over competitors gives \eqref{eq:central-lower}, without exchanging an infimum and a supremum. Nonnegative loss gives the zero truncation and the stated case $t_N=0$.

For the upper assume \eqref{eq:beta}. Choose an attained optimizer $(C,h)$ at width $m_+(M)$ and cap $\beta_N$. Let $r_C(x)$ be the first nearest point to $f(x)$ in a fixed ordering of $C$. After each attempt of $\mathcal G$, accept $x$ with probability $h(x)/\beta_N$, export $T_Jr_C(x)$ and commit; otherwise return to the start routing label. All these operations are Borel. The choice $J$ is used within the current attempt; any part retained thereafter is charged through the live role or exported decoder, not a free seed. The exports are assumed jointly measurable and take at most $D_0|C|$ distinct values, so their finite decoder and transitions are Borel as well.

Conditional on every surviving attempt history, the acceptance probability is at least
\begin{equation}\label{eq:hazard}
 p=\frac{l_+}{\beta_N}-a_+=\frac1{b_+}>0.
\end{equation}
Consequently the probability of reaching attempt $j+1$ is at most $(1-p)^j$, by induction on conditional survival probabilities. In particular commitment is almost sure and the expected number of attempts is at most $1/p=b_+$. Each begun attempt has conditional expected cost at most $\kappa_+$, not merely a cost conditional on a favorable outcome. Summing the predictable attempt costs gives actual expected cost at most $\kappa_+b_+=N$.

The conditional unnormalized loss contributed by an attempt is at most $u_+\Q_{m_+(M)}(\beta_N)/\beta_N+e_+$. Sum its contributions on the disjoint terminal events, using the same survival bound. The result is at most
\[
 b_+\left[\frac{u_+}{\beta_N}\Q_{m_+(M)}(\beta_N)+e_+\right],
\]
which is \eqref{eq:central-upper}. The program has $B$ live roles and at most $D_0m_+(M)$ terminal roles. Lemma~\ref{lem:roles} gives exactly the sufficient widths in \eqref{eq:widths}. A checkpoint counts only the terminal width. No physical renewal law was used: the conditional hazard controls survival even if the hidden state after rejection differs arbitrarily from its former state.

If $\beta_N<1$, commit after the first attempt, and use an optimizer at cap one. The loss certificate with $q=1$ gives $u_+\Q_{m_+(M)}(1)+e_+$, and the whole-attempt cost is at most $\kappa_+$. This proves all remaining assertions.
\end{proof}

\section{Scope and further mathematical questions}\label{sec:r73-scope}
The normal form is a classification of exact finite stochastic
\emph{input--response realization} on the compactly reversible observable
quotient. It applies to every finite prepared unitary command family, not
only a unary rotation or a fixed protected dimension. Its size formula
counts a prescribed initial live label and bounded terminal probabilities.
It neither exposes a simulator's private state to the experimenter nor
optimizes an experimenter who may choose a less demanding command policy.
General several-command boundary-optimal feedback remains a separate
problem.

The algebraic synthesis theorem takes an exact supplied finite linear
representation. Probability validity and compact reversibility are part
of this input class; both follow from the prepared unitary model in its
corollary. No algorithm for learning an unknown predictive representation
from finite observations is inferred. The returned label minimum is
exact, but the elementary group bound and real-algebraic search can be
very large. Transition and decoder descriptions, sampling, arithmetic
precision, program length and transient working storage remain additional
resources.

The common-background damping plateau is a proved positive interval of
unchanged exact memory. Algebraic data make a valid cutoff computable;
they do not make the first transition point a closed-form expression.
The periodic partition and leakage laws retain their declared spectral
support and probability-valid mixture. The arbitrary-phase arithmetic
envelope retains both logarithmic exponents and its badly-approximable
criterion, but no particular irrationality exponent is deduced for the
algebraic qubit angle with $e^{i\vartheta}=(3+4i)/5$.

The full-history optimum, exact finite-memory attainment, finite-memory
density, approximation risk at fixed memory, and exact response
simulation are different assertions. The previously proved robust tied
criterion on the rigid herald class remains in force. The older general
$M$ squared-sign risk curve, intermediate phase-type timing frontier,
sharp arithmetic leading constants and matching full computational
resource lower bounds are not settled by the new normal form.

The actual-history transfer theorem above and the performed-law,
variational, dual, spectral, reset, discrimination and geometric theorems
in the current supplement retain their complete proofs. The independently
compiled absolute-calibration companion and the original ordered
measurement and two-stage reset realizations retain their own source
identities. Results requiring additional A1, A2 or broader Theta
construction hypotheses are not consequences of a finite-group response
classification. The general research problem remains the relation
between information acquired by causal experiments and the resources
required to realize their optimal use.
\begin{thebibliography}{99}
\bibitem{Blackwell} D. Blackwell, \emph{Equivalent comparisons of experiments}, Ann. Math. Statist. \textbf{24} (1953), 265--272. DOI: 10.1214/aoms/1177729032.
\bibitem{CGM97} J. A. Cuesta-Albertos, A. Gordaliza and C. Matr\'an, \emph{Trimmed $k$-means: an attempt to robustify quantizers}, Ann. Statist. \textbf{25} (1997), 553--576. DOI: 10.1214/aos/1031833664.
\bibitem{GG99} L. A. Garc\'ia-Escudero and A. Gordaliza, \emph{Robustness properties of $k$ means and trimmed $k$ means}, J. Amer. Statist. Assoc. \textbf{94} (1999), 956--969. DOI: 10.1080/01621459.1999.10474200.
\bibitem{SC} C. R. Shalizi and J. P. Crutchfield, \emph{Computational mechanics: pattern and prediction, structure and simplicity}, J. Statist. Phys. \textbf{104} (2001), 817--879. DOI: 10.1023/A:1010388907793.
\bibitem{BC} N. Barnett and J. P. Crutchfield, \emph{Computational mechanics of input--output processes: structured transformations and the $\varepsilon$-transducer}, J. Statist. Phys. \textbf{161} (2015), 404--451; arXiv:1412.2690.
\bibitem{AMTW} A. Ambainis, M. Mosca, A. Tapp and R. de Wolf, \emph{Private quantum channels}, Proc. 41st Annual Symposium on Foundations of Computer Science (2000), 547--553. DOI: 10.1109/SFCS.2000.892142.
\bibitem{Buscemi} F. Buscemi, \emph{Comparison of quantum statistical models: equivalent conditions for sufficiency}, Comm. Math. Phys. \textbf{310} (2012), 625--647. DOI: 10.1007/s00220-012-1421-3.
\bibitem{BR14} N. Bertschinger and J. Rauh, \emph{The Blackwell relation defines no lattice}, Proc. IEEE International Symposium on Information Theory (2014), 2479--2483. DOI: 10.1109/ISIT.2014.6875280; arXiv:1401.3146.
\bibitem{Banerjee05} A. Banerjee, X. Guo and H. Wang, \emph{On the optimality of conditional expectation as a Bregman predictor}, IEEE Trans. Inform. Theory \textbf{51} (2005), 2664--2669. DOI: 10.1109/TIT.2005.850145.
\bibitem{Stengel96} B. von Stengel, \emph{Efficient computation of behavior strategies}, Games Econom. Behav. \textbf{14} (1996), 220--246. DOI: 10.1006/game.1996.0050.
\bibitem{Balder89} E. J. Balder, \emph{On compactness of the space of policies in stochastic dynamic programming}, Stochastic Process. Appl. \textbf{32} (1989), 141--150. DOI: 10.1016/0304-4149(89)90058-6.
\bibitem{Barry14} J. Barry, D. T. Barry and S. Aaronson, \emph{Quantum partially observable Markov decision processes}, Phys. Rev. A \textbf{90} (2014), 032311. DOI: 10.1103/PhysRevA.90.032311; arXiv:1406.2858.
\bibitem{Ying18} S. Ying and M. Ying, \emph{Reachability analysis of quantum Markov decision processes}, Inform. Comput. \textbf{263} (2018), 31--51. DOI: 10.1016/j.ic.2018.09.001; arXiv:1406.6146.
\bibitem{Saldi25} N. Saldi, S. Sanjari and S. Y\"uksel, \emph{Quantum Markov decision processes: general theory, approximations, and classes of policies}, SIAM J. Control Optim. \textbf{63} (2025), 1156--1183. DOI: 10.1137/24M1643864; arXiv:2402.14649.
\bibitem{Seidenberg54} A. Seidenberg, \emph{A new decision method for elementary algebra}, Ann. of Math. (2) \textbf{60} (1954), 365--374. DOI: 10.2307/1969640.
\bibitem{Chatterjee15} K. Chatterjee, M. Chmel\'ik, R. Gupta and A. Kanodia, \emph{Optimal cost almost-sure reachability in POMDPs}, Proc. AAAI Conference on Artificial Intelligence \textbf{29} (2015). DOI: 10.1609/aaai.v29i1.9683.
\bibitem{Andriushchenko23} R. Andriushchenko, A. Bork, M. \v{C}e\v{s}ka, S. Junges, J.-P. Katoen and M. Mac\'ak, \emph{Search and explore: symbiotic policy synthesis in POMDPs}, Computer Aided Verification (2023); arXiv:2305.14149v2.
\bibitem{Halmos69} P. R. Halmos, \emph{Two subspaces}, Trans. Amer. Math. Soc. \textbf{144} (1969), 381--389. DOI: 10.1090/S0002-9947-1969-0251519-5.
\bibitem{Kessler21} D. A. Kessler, E. Barkai and K. Ziegler, \emph{First-detection time of a quantum state under random probing}, Phys. Rev. A \textbf{103} (2021), 022222. DOI: 10.1103/PhysRevA.103.022222; arXiv:2012.01763.
\bibitem{DelVecchio26} G. Del Vecchio Del Vecchio and S. N. Majumdar, \emph{Optimal detection of quantum states via projective measurements}, J. Phys. A: Math. Theor. \textbf{59} (2026), 035001. DOI: 10.1088/1751-8121/ae34fe; arXiv:2509.08556.
\bibitem{Companion} Q. Qi, \emph{General Theta Foundations I: Acquired Geometry and Causal Resource Transfer}, complete technical companion on absolute calibration, exact stopping resources and nonintervening sensing, 2026. The complete retained R58 native source and its independently built companion accompany this manuscript; their identities are fixed in the source-provenance and preservation records.
\bibitem{Tzeng92} W.-G. Tzeng, \emph{A polynomial-time algorithm for the equivalence of probabilistic automata}, SIAM J. Comput. \textbf{21} (1992), 216--227. DOI: 10.1137/0221017.
\bibitem{LiYuYing12} Y. Li, N. Yu and M. Ying, \emph{Termination of nondeterministic quantum programs}, arXiv:1201.0891 (2012).
\bibitem{Fu24} J. Fu, H. Jiang, M. Xu, Y. Deng and Z.-B. Li, \emph{Algorithmic analysis of termination problems for nondeterministic quantum programs}, arXiv:2402.15827 (2024).

\bibitem{LiuVerdu18} J. Liu and S. Verd\'u, \emph{Rejection sampling and noncausal sampling under moment constraints}, Proc. IEEE International Symposium on Information Theory (2018), 1565--1569. DOI: 10.1109/ISIT.2018.8437857.
\bibitem{GocFlamich24} D. Goc and G. Flamich, \emph{On channel simulation with causal rejection samplers}, Proc. IEEE International Symposium on Information Theory (2024). DOI: 10.1109/ISIT57864.2024.10619339; examined text arXiv:2401.16579v2 (3 May 2024).
\bibitem{LiLiu26} Z. Li and T. Liu, \emph{Ranking experiments under sequential sampling}, arXiv:2608.19897v2 (28 August 2026). Version 1 appeared under the title \emph{The order of binary experiments under endogenous stopping}, with Z. Li as sole author; the version-2 statement is the one compared here.
\bibitem{Ying14pre} S. Ying and M. Ying, \emph{Reachability analysis of quantum Markov decision processes}, arXiv:1406.6146 (2014), Theorem 3.6; published journal version listed above as \cite{Ying18}.

\bibitem{KaraYuksel20} A. D. Kara and S. Y\"uksel, \emph{Robustness to incorrect system models in stochastic control}, SIAM J. Control Optim. \textbf{58} (2020), 1144--1182. DOI: 10.1137/18M1208058; examined text arXiv:1803.06046v7.
\bibitem{Demirci25} Y. E. Demirci, A. D. Kara and S. Y\"uksel, \emph{Sensitivity of filter kernels and robustness bounds to transition and measurement kernel perturbations in partially observable stochastic control}, arXiv:2508.10658v3 (2025).
\bibitem{Watrous18} J. Watrous, \emph{The Theory of Quantum Information}, Cambridge University Press, 2018, Chapter 3, Corollary 3.40.
\bibitem{KRN65} K. Kuratowski and C. Ryll-Nardzewski, \emph{A general theorem on selectors}, Bull. Acad. Polon. Sci. S\'er. Sci. Math. Astronom. Phys. \textbf{13} (1965), 397--403.

\bibitem{Charnes63} A. Charnes, W. W. Cooper and K. O. Kortanek, \emph{Duality in semi-infinite programs and some works of Haar and Carath\'eodory}, Management Sci. \textbf{9} (1963), 209--228. DOI: 10.1287/mnsc.9.2.209.
\bibitem{Shapiro09} A. Shapiro, \emph{Semi-infinite programming, duality, discretization and optimality conditions}, Optimization \textbf{58} (2009), 133--161. DOI: 10.1080/02331930902730070; author preprint, Optimization Online (2008).
\bibitem{Dufour25} F. Dufour and T. Prieto-Rumeau, \emph{Absorbing Markov decision processes and their occupation measures}, SIAM J. Control Optim. \textbf{63} (2025), 676--698. DOI: 10.1137/24M1661959. Related author preprint examined: \emph{A note on weak compactness of occupation measures for an absorbing Markov decision process}, arXiv:2402.10672v2 (2024).
\bibitem{Franc23} V. Franc, D. Prusa and V. Voracek, \emph{Optimal strategies for reject option classifiers}, J. Mach. Learn. Res. \textbf{24} (2023), no. 11, 1--49.
\bibitem{Jencova16} A. Jen\v{c}ov\'a, \emph{Comparison of quantum channels and statistical experiments}, arXiv:1512.07016v2 (2016). The separately indexed shorter version arXiv:1601.06370 was withdrawn in favor of this extended text.
\bibitem{LiElGamal18} C. T. Li and A. El Gamal, \emph{Strong functional representation lemma and applications to coding theorems}, IEEE Trans. Inform. Theory \textbf{64} (2018), 6967--6978; arXiv:1701.02827v4.
\bibitem{AldousShepp87} D. Aldous and L. Shepp, \emph{The least variable phase type distribution is Erlang}, Comm. Statist. Stochastic Models \textbf{3} (1987), 467--473. DOI: 10.1080/15326348708807067.
\bibitem{Sion58} M. Sion, \emph{On general minimax theorems}, Pacific J. Math. \textbf{8} (1958), 171--176. DOI: 10.2140/pjm.1958.8.171.
\bibitem{PiunovskiyZhang24} A. Piunovskiy and Y. Zhang, \emph{Extreme occupation measures in Markov decision processes with an absorbing state}, SIAM J. Control Optim. \textbf{62} (2024), 65--90. DOI: 10.1137/23M1572398. Author text examined: \emph{Extreme occupation measures in Markov decision processes with a cemetery}, arXiv:2307.03158v1.
\bibitem{DufourMixtures} F. Dufour and T. Prieto-Rumeau, \emph{Absorbing Markov decision processes: geometric properties and sufficiency of finite mixtures of deterministic policies}, arXiv:2512.17511; examined HTML at version-1 URL, Theorems 4.3--4.4 and Proposition 4.5.
\bibitem{Mereghetti24} C. Mereghetti, B. Palano and P. Raucci, \emph{Unary quantum finite state automata with control language}, Appl. Sci. \textbf{14} (2024), no. 4, 1490. DOI: 10.3390/app14041490.
\bibitem{RabinScott59} M. O. Rabin and D. Scott, \emph{Finite automata and their decision problems}, IBM J. Res. Develop. \textbf{3} (1959), 114--125. DOI: 10.1147/rd.32.0114.
\bibitem{Jaeger00} H. Jaeger, \emph{Observable operator models for discrete stochastic time series}, Neural Comput. \textbf{12} (2000), 1371--1398. DOI: 10.1162/089976600300015411.
\bibitem{ThonJaeger15} M. Thon and H. Jaeger, \emph{Links between multiplicity automata, observable operator models and predictive state representations---a unified learning framework}, J. Mach. Learn. Res. \textbf{16} (2015), 103--147.
\bibitem{MonrasWinter14} A. Monras and A. Winter, \emph{Quantum learning of classical stochastic processes: the completely-positive realization problem}, LIPIcs \textbf{27}, TQC 2014 (2014), 99--109. DOI: 10.4230/LIPIcs.TQC.2014.99.

\bibitem{Heller65} A. Heller, \emph{On stochastic processes derived from Markov chains}, Ann. Math. Statist. \textbf{36} (1965), 1286--1291. DOI: 10.1214/aoms/1177700000.
\bibitem{CarlylePaz71} J. W. Carlyle and A. Paz, \emph{Realizations by stochastic finite automata}, J. Comput. System Sci. \textbf{5} (1971), 26--40. DOI: 10.1016/S0022-0000(71)80005-3.
\bibitem{Fanizza24} M. Fanizza, J. Lumbreras and A. Winter, \emph{Quantum theory in finite dimension cannot explain every general process with finite memory}, Comm. Math. Phys. \textbf{405} (2024), article 50. DOI: 10.1007/s00220-023-04913-4; arXiv:2209.11225. In particular see Theorem 2 and Section 4.2 for the noise-dependent worst-phase classical-memory lower bound.
\bibitem{DmitrievDynkin46} N. Dmitriev and E. Dynkin, \emph{On characteristic roots of stochastic matrices}, Izv. Akad. Nauk SSSR Ser. Mat. \textbf{10} (1946), no. 2, 167--184. Math-Net.Ru: im3595 (Russian).
\bibitem{JohnsonPaparella17} C. R. Johnson and P. Paparella, \emph{A matricial view of the Karpelevi\v c theorem}, Linear Algebra Appl. \textbf{520} (2017), 1--15. DOI: 10.1016/j.laa.2017.01.009; arXiv:1611.06970v2.
\bibitem{KirklandSmigoc22} S. Kirkland and H. \v Smigoc, \emph{Stochastic matrices realising the boundary of the Karpelevi\v c region}, Linear Algebra Appl. \textbf{635} (2022), 116--138. DOI: 10.1016/j.laa.2021.11.016; arXiv:2110.01040v1.
\bibitem{BenvenutiFarina04} L. Benvenuti and L. Farina, \emph{A tutorial on the positive realization problem}, IEEE Trans. Automat. Control \textbf{49} (2004), 651--664. DOI: 10.1109/TAC.2004.826715.
\bibitem{BenvenutiFarina99} L. Benvenuti and L. Farina, \emph{An example of how positivity may force realizations of ``large'' dimension}, Systems Control Lett. \textbf{36} (1999), 261--266. DOI: 10.1016/S0167-6911(98)00098-X.
\bibitem{OCinneide90} C. A. O'Cinneide, \emph{Characterization of phase-type distributions}, Comm. Statist. Stochastic Models \textbf{6} (1990), 1--57. DOI: 10.1080/15326349908807134.
\bibitem{Maier91} R. S. Maier, \emph{The algebraic construction of phase-type distributions}, Comm. Statist. Stochastic Models \textbf{7} (1991), 573--602. DOI: 10.1080/15326349108807207.
\bibitem{HoneVarona19} A. N. W. Hone and J. L. Varona, \emph{Continued fractions and irrationality exponents for modified Engel and Pierce series}, Monatsh. Math. \textbf{190} (2019), 501--516. DOI: 10.1007/s00605-018-1244-1. Author version arXiv:1810.10127v1, equation (3.3).
\bibitem{Nussbaum98} R. D. Nussbaum and M. Scheutzow, \emph{Admissible arrays and a nonlinear generalization of Perron--Frobenius theory}, J. London Math. Soc. (2) \textbf{58} (1998), 526--544. DOI: 10.1112/S0024610798006796.
\bibitem{Brown64} D.~R. Brown, On clans of non-negative matrices,
\emph{Proc. Amer. Math. Soc.} \textbf{15} (1964), 671--674,
doi:10.1090/S0002-9939-1964-0165504-6.
\bibitem{Schwarz64} \v{S}. Schwarz, On the structure of the semigroup of
stochastic matrices, \emph{Publ. Math. Inst. Hungar. Acad. Sci.}
\textbf{9} (1964), 297--311, especially Theorems~7--8.
\bibitem{Flor69} P. Flor, On groups of non-negative matrices,
\emph{Compositio Math.} \textbf{21} (1969), 376--382,
\url{https://www.numdam.org/article/CM_1969__21_4_376_0.pdf}.
\bibitem{Serre07} J.-P. Serre, Bounds for the orders of the finite subgroups
of $G(k)$, in \emph{Group Representation Theory}, M. Geck, D. Testerman
and J. Th\'evenaz (eds.), EPFL Press, Lausanne, 2007, 405--450;
author lecture text, April 2006, archived as arXiv:1011.0346.
\bibitem{BabaiBealsRockmore93} L. Babai, R. Beals and D. Rockmore,
Deciding finiteness of matrix groups in deterministic polynomial time,
\emph{Proceedings of ISSAC 1993}, 117--126,
doi:10.1145/164081.164104.
\bibitem{KieferWachter14} S. Kiefer and B. Wachter, Stability and complexity
of minimising probabilistic automata, in \emph{ICALP 2014, Part II},
Lecture Notes in Computer Science \textbf{8573}, 268--279,
doi:10.1007/978-3-662-43951-7\_23; arXiv:1404.6673, Section~4.2.
\bibitem{Schwarz62} \v{S}. Schwarz, Homomorphisms of a completely simple
semigroup onto a group, \emph{Matematicko-fyzik\'alny \v{c}asopis}
\textbf{12} (1962), no.~4, 293--300;
\url{https://www.dml.cz/handle/10338.dmlcz/126333}.
\end{thebibliography}

\end{document}
